\documentclass[twoside,11pt]{article}

\usepackage[utf8]{inputenc}
\usepackage[T2A]{fontenc}
\usepackage[english,russian]{babel}
\usepackage{amsmath,amsfonts,amssymb}
\usepackage{graphicx}
\usepackage{multirow}
\usepackage{booktabs}
\usepackage{tabularx}
\usepackage{longtable}
\usepackage{array}
\usepackage{url}
\usepackage{hyperref}
\usepackage{enumitem}
\usepackage{placeins}
\usepackage{float}
\usepackage{xcolor}
\usepackage{caption}
\usepackage{geometry}

\geometry{
  a4paper,
  left=25mm,
  right=20mm,
  top=22mm,
  bottom=24mm
}

\hypersetup{
  colorlinks=true,
  linkcolor=blue,
  citecolor=blue,
  urlcolor=blue,
  pdftitle={Curriculum-KAG: от результатов обучения к проверяемому учебному плану},
  pdfauthor={Murat Kozhanov}
}

\renewcommand{\topfraction}{0.95}
\renewcommand{\bottomfraction}{0.95}
\renewcommand{\textfraction}{0.05}
\renewcommand{\floatpagefraction}{0.90}
\setcounter{topnumber}{4}
\setcounter{bottomnumber}{4}
\setcounter{totalnumber}{6}

\newcommand{\R}{\mathbb{R}}
\newcommand{\Ccal}{\mathcal{C}}
\newcommand{\Ocal}{\mathcal{O}}
\newcommand{\Pcal}{\mathcal{P}}
\newcommand{\Dcal}{\mathcal{D}}
\newcommand{\Ecal}{\mathcal{E}}
\newcommand{\Gcal}{\mathcal{G}}
\newcommand{\clip}{\operatorname{clip}}
\newcommand{\sigmoid}{\operatorname{sigmoid}}
\newcommand{\placeholderfigure}[2]{%
  \fbox{\parbox[c][5.2cm][c]{0.92\linewidth}{%
    \centering
    \textbf{Место для рисунка:} \texttt{#1}\\[6pt]
    #2
  }}%
}
\newcommand{\systemname}{\textsc{Curriculum-KAG}}
\newenvironment{plainbox}{%
  \begin{quote}\small\noindent\textbf{Простое объяснение.}\ }{%
  \end{quote}}

\begin{document}

\begin{center}
{\Large\bf Curriculum-KAG: от результатов обучения к проверяемому учебному плану на основе экспертных данных ЕПВО и формальной оптимизации}\\[8pt]

{\large\bf Murat Kozhanov}\\[3pt]
Abylkas Saginov Karaganda Technical University, Karaganda, Kazakhstan\\
Doctoral School of Applied Informatics and Applied Mathematics, Obuda University, Budapest, Hungary\\
\texttt{m.kozhanov@ktu.edu.kz}\\[8pt]
\end{center}

\noindent\textbf{Аннотация.}
Автоматизированное проектирование образовательной программы требует одновременно обеспечить содержательную связь дисциплин с результатами обучения, корректную последовательность пререквизитов, нормативный объём кредитов, баланс семестров, принадлежность дисциплин выбранному направлению подготовки и объяснимость решений для эксперта. Генерация только языковой моделью не гарантирует соблюдение этих условий и затрудняет воспроизводимость. В статье представлена система \systemname{}, объединяющая нормализованный репозиторий Единой платформы высшего образования Республики Казахстан (ЕПВО), гибридный семантико-лексический поиск, граф дисциплин и пререквизитов, экспертные оценки связей ``дисциплина--результат обучения'', вероятностную оценку покрытия и многокритериальную оптимизацию учебного плана.

Из реестра сформирован воспроизводимый набор данных, включающий 12\,215 образовательных программ, 408\,638 исходных записей дисциплин, 124\,521 исходный результат обучения, 191\,292 нормализованные дисциплины и 932\,483 нормализованные экспертные связи. Разбиение train/validation/test выполнено на уровне программ с фиксированным seed~42, что снижает риск утечки одинакового контекста одной программы между выборками. На независимой тестовой выборке из 6\,000 пар дообученная модель \texttt{epvo-sbert-finetuned-40k} достигла ROC-AUC 0,7666, PR-AUC 0,7681 и F1 0,7188 против 0,7031, 0,6946 и 0,6886 у исходной multilingual SBERT. Cached-модель 12k улучшила классификацию до ROC-AUC 0,7748, PR-AUC 0,7746 и F1 0,7268, однако сохранена как кандидат, а не автоматически внедрена вместо проверенного рабочего baseline. Отдельная модель с Multiple Negatives Ranking Loss улучшила ранжирование до Recall@10 0,7071, MRR 0,7815 и nDCG@10 0,6923. Контролируемые пилоты GNN и LSTM не превзошли SBERT, что позволило избежать необоснованного внедрения более сложных архитектур. В контрольной регрессии девять вариантов A/B/C для трёх междисциплинарных бакалаврских программ получили целевые 240 кредитов, полное реальное покрытие LO и ноль жёстких нарушений. Дополнительный контроль магистратуры подтвердил A/B/C по 120 кредитов без жёстких нарушений, а независимый dry-run докторантуры D094 сформировал 180 кредитов, пять реальных профильных дисциплин, ноль bridge-модулей, 100\% покрытия LO и допустимую нагрузку $30/30/30/30/27/33$.

Научная новизна состоит в совместном использовании национального экспертного корпуса, двухконтурной модели ``классификация связи + ранжирование'', вероятностного агрегирования покрытия, иерархического отбора дисциплин по ЕПВО и constrained-синтеза плана с проверяемыми ограничениями. Система сохраняет происхождение данных, показывает причину выбора дисциплины и включает экспертную обратную связь, поэтому рассматривается как инструмент поддержки принятия решений, а не как замена академического эксперта.

\medskip
\noindent\textbf{Ключевые слова:}
образовательная программа, результаты обучения, ЕПВО, KAG, SBERT, семантический поиск, граф знаний, NSGA-II, пререквизиты, многокритериальная оптимизация, explainable AI, human-in-the-loop.

\bigskip
\noindent\textbf{English title.}
Curriculum-KAG: From Learning Outcomes to a Verifiable Curriculum Using EPVO Expert Data and Constrained Optimization

\medskip
\noindent\textbf{Abstract.}
Automated curriculum design must simultaneously align courses with programme learning outcomes, preserve prerequisite order, satisfy credit and semester-load regulations, respect selected educational domains, and provide expert-readable explanations. Pure language-model generation does not guarantee these properties. This paper presents \systemname{}, an implemented decision-support system that combines a normalized repository of Kazakhstan's Unified Higher Education Platform (EPVO), hybrid semantic and lexical retrieval, a course--prerequisite graph, expert course--learning-outcome judgements, probabilistic coverage aggregation, and constrained multi-objective curriculum optimization.

The reproducible dataset contains 12,215 programmes, 408,638 raw course records, 124,521 raw learning outcomes, 191,292 normalized courses, and 932,483 normalized expert links. A programme-level frozen split with seed 42 is used to limit contextual leakage. On 6,000 independent test pairs, the fine-tuned \texttt{epvo-sbert-finetuned-40k} model achieved ROC-AUC 0.7666, PR-AUC 0.7681, and F1 0.7188, compared with 0.7031, 0.6946, and 0.6886 for the base multilingual SBERT. A cached 12k candidate improved classification to ROC-AUC 0.7748, PR-AUC 0.7746, and F1 0.7268, but was retained as an experimental candidate rather than silently replacing the validated operational baseline. A separate Multiple Negatives Ranking Loss pilot improved retrieval to Recall@10 0.7071, MRR 0.7815, and nDCG@10 0.6923. Controlled GNN and LSTM pilots did not outperform SBERT and were therefore not promoted to the operational pipeline. Nine A/B/C variants from three interdisciplinary bachelor's programmes reached 240 credits, complete real-LO coverage, and zero hard violations. A master's control passed all A/B/C checks at 120 credits; an independent D094 doctoral dry-run produced 180 credits, five real professional courses, no bridge modules, 100\% LO coverage, and admissible semester loads. The contribution is a reproducible expert-data pipeline and a two-stage, constraint-aware, explainable approach to curriculum synthesis.

\medskip
\noindent\textbf{Keywords:}
curriculum design, learning outcomes, EPVO, KAG, SBERT, dense retrieval, knowledge graph, NSGA-II, prerequisite graph, explainable AI, human-in-the-loop.

\section{Введение}

Современная образовательная программа должна проектироваться от ожидаемых результатов обучения к содержанию дисциплин, учебным активностям и оцениванию. Такая логика соответствует constructive alignment \cite{Biggs1996} и требованиям outcome-based education: результат обучения должен быть не декларацией, а проверяемым утверждением о знаниях, умениях и профессиональном поведении выпускника. Международные рамки качества также требуют явной связи результатов обучения с учебным планом и систематического использования данных оценивания для непрерывного улучшения программы \cite{ABET2025,CDIO2022,Tuning2003}.

На практике разработчик программы сталкивается с несколькими взаимозависимыми задачами:
\begin{enumerate}[label=(\arabic*)]
  \item выбрать дисциплины, действительно релевантные целям и результатам обучения;
  \item не допустить случайных или чужих предметных областей;
  \item расположить дисциплины так, чтобы базовые знания предшествовали продвинутым;
  \item обеспечить заданный объём кредитов и допустимую нагрузку каждого семестра;
  \item сформировать разные, но одинаково корректные альтернативы плана;
  \item показать эксперту, почему выбрана каждая дисциплина и какой результат она поддерживает;
  \item сохранить возможность ручного подтверждения, ослабления или отклонения предложенной связи.
\end{enumerate}

Языковая модель способна предложить названия дисциплин и текстовые описания, но сама по себе не является решателем ограничений. Она может повторять дисциплины, ставить преддипломную практику в первый семестр, связывать общеобразовательный курс с профессиональным результатом без доказательства или нарушать кредитную структуру. Следовательно, задача должна формулироваться не как свободная генерация текста, а как поиск и оптимизация в пространстве допустимых учебных планов.

Retrieval-Augmented Generation (RAG) связывает параметрическую модель с внешней базой знаний \cite{Lewis2020}; плотный поиск использует представления текстов в общем векторном пространстве \cite{Karpukhin2020}. Knowledge-Augmented Generation (KAG) расширяет эту идею структурированным знанием, отношениями и правилами предметной области \cite{Liang2024,Hogan2021}. В \systemname{} принцип KAG переносится с генерации ответа на проектирование целой образовательной программы: извлечённые дисциплины образуют кандидатов, граф задаёт отношения, а оптимизатор выбирает и распределяет элементы под формальными ограничениями.

Основой эмпирической части стал Реестр образовательных программ ЕПВО. Реестр выполняет учётно-информационную функцию и содержит сведения об образовательных программах организаций высшего и послевузовского образования Казахстана \cite{EPVORegistry}; правила его ведения определены нормативным актом Республики Казахстан \cite{AdiletRegistry}. Карточки ЕПВО важны не только как каталог названий. Они содержат направления и группы образовательных программ, цели, результаты обучения, дисциплины, их описания на трёх языках и рассеянные по карточкам экспертные оценки связей дисциплин с результатами обучения. Эти оценки позволяют обучать модель не на искусственно придуманных соответствиях, а на реальной экспертной практике.

В статье решаются следующие исследовательские вопросы:
\begin{itemize}
  \item \textbf{RQ1:} улучшает ли предметное дообучение multilingual SBERT качество определения связи ``дисциплина--результат обучения''?
  \item \textbf{RQ2:} улучшают ли hard-negative mining и контрастное ranking-обучение не только классификационные, но и поисковые метрики?
  \item \textbf{RQ3:} дают ли GNN и LSTM преимущество над сильным текстовым baseline при текущем представлении данных?
  \item \textbf{RQ4:} возможно ли превратить прогноз связей в корректный план, удовлетворяющий кредитам, пререквизитам, семестровой нагрузке и предметным квотам?
  \item \textbf{RQ5:} можно ли сохранить объяснимость, происхождение данных и роль эксперта в полном цикле генерации?
\end{itemize}

Основные результаты работы:
\begin{enumerate}
  \item разработан многоуровневый конвейер \texttt{raw EPVO $\rightarrow$ normalized EPVO $\rightarrow$ approved repository $\rightarrow$ generation candidates};
  \item собран воспроизводимый корпус из 935\,151 пар с программно-уровневым разбиением и контрольными суммами;
  \item дообучена SBERT-модель, улучшившая ROC-AUC на 0,0635 и PR-AUC на 0,0735 относительно исходной модели;
  \item экспериментально показано, что ranking-loss улучшает порядок кандидатов, хотя не становится лучшим классификатором;
  \item реализован детерминированный constrained-планировщик A/B/C с repair-эвристиками, точным кредитным балансированием и независимой жёсткой верификацией; NSGA-II сохранён как экспериментальная ветка для программ без явного EPVO-скоупа и не заявляется как основной производственный оптимизатор;
  \item реализованы объяснения, bridge-модули, замена bridge реальными дисциплинами ЕПВО и экспертная обратная связь;
  \item создан научный контур Dataset Passport и Model Benchmark, фиксирующий данные, split, seed, checksum и метрики.
\end{enumerate}

\subsection*{Как система работает без технических терминов}

Разработчик сначала указывает, кого и чему должна научить программа, её уровень
образования, направление ЕПВО и нормативный объём кредитов. Далее система не
``придумывает весь план'', а последовательно выполняет семь проверяемых действий:
\begin{enumerate}
  \item оставляет дисциплины нужного уровня и выбранной группы ЕПВО;
  \item сравнивает описание каждой дисциплины с каждым результатом обучения;
  \item учитывает, какие связи ранее были отмечены экспертами ЕПВО;
  \item строит граф ``что изучается раньше и что открывает дальше'';
  \item собирает несколько вариантов A/B/C под заданные кредиты и квоты;
  \item распределяет дисциплины по семестрам и проверяет пререквизиты;
  \item показывает человеку причины выбора и не сохраняет новый набор планов,
        если хотя бы один вариант нарушает жёсткие правила.
\end{enumerate}

\begin{plainbox}
SBERT здесь похож на помощника библиотекаря: он находит содержательно близкие
дисциплины, даже если слова не совпадают буквально. Оптимизатор похож на
составителя расписания: из найденных вариантов он собирает целый план. Наконец,
верификатор похож на независимого контролёра: он заново проверяет кредиты,
уровень, последовательность и покрытие результатов. Ни один из этих трёх
компонентов по отдельности не заменяет остальные.
\end{plainbox}

\begin{figure}[H]
\centering
\placeholderfigure{fig01\_curriculum\_kag\_architecture.pdf}{Концептуальная архитектура: мастер ОП $\rightarrow$ scoped EPVO repository $\rightarrow$ hybrid retrieval $\rightarrow$ course--LO scoring $\rightarrow$ knowledge graph $\rightarrow$ deterministic constrained planner $\rightarrow$ verifier $\rightarrow$ explainability and expert feedback. NSGA-II показан отдельной экспериментальной веткой.}
\caption{Общая архитектура \systemname{}: от scoped-данных ЕПВО до проверенного плана и экспертной обратной связи.}
\label{fig:architecture}
\end{figure}

\section{Связанные исследования}

\subsection{Outcome-based design и непрерывное улучшение}

Constructive alignment связывает ожидаемые результаты, учебную деятельность и оценивание \cite{Biggs1996}. В логике ABET результаты студентов должны быть документированы, оценены, а результаты оценки систематически использованы для улучшения программы \cite{ABET2025}. Стандарт CDIO~3.0 требует интегрированной учебной программы, в которой дисциплины взаимно поддерживают профессиональные, личностные и межличностные компетенции \cite{CDIO2022}. Методология Tuning использует результаты обучения и компетенции как точки сопоставления программ при сохранении институциональной гибкости \cite{Tuning2003}. Эти подходы обосновывают матрицу ``дисциплина--LO'', но не дают готового алгоритма автоматического синтеза полного плана.

\subsection{Семантические представления и гибридный поиск}

BERT показал эффективность двунаправленного предобучения трансформеров \cite{Devlin2019}. SBERT преобразует предложения в векторы, допускающие быстрый поиск по косинусной близости \cite{Reimers2019}; multilingual distillation позволяет поддерживать многоязычное пространство предложений \cite{Reimers2020}. Использованная базовая модель основана на MPNet, совмещающем masked и permuted language modeling \cite{Song2020}. Для образовательных данных это существенно, поскольку названия, описания и результаты обучения в ЕПВО представлены на русском, казахском и английском языках.

Плотный поиск эффективен для смысловой близости, но может пропустить точные термины. BM25, напротив, хорошо использует лексические совпадения \cite{Robertson2009}. Поэтому в системе применяется гибридный score. Аналогичный принцип сочетания параметрической и непараметрической памяти лежит в основе RAG \cite{Lewis2020}, а KAG добавляет структурированные отношения и знания предметной области \cite{Liang2024}.

\subsection{Графы и последовательности учебного плана}

Граф знаний представляет сущности и отношения в единой структуре \cite{Hogan2021}. Учебный план естественно моделируется ориентированным графом, где вершины --- дисциплины, а дуги --- пререквизиты. Подобная модель используется в curricular analytics для исследования сложности, задерживающих и блокирующих дисциплин \cite{Heileman2018}.

Graph Convolutional Networks агрегируют признаки соседей графа \cite{Kipf2017}, а GraphSAGE обучает индуктивную функцию агрегации для новых вершин \cite{Hamilton2017}. LSTM предназначена для моделирования длительных зависимостей в последовательностях \cite{Hochreiter1997}. Эти архитектуры выглядят естественными кандидатами для графа дисциплин и последовательности семестров. Однако сложность модели сама по себе не гарантирует улучшение: необходимы одинаковые split, признаки, отрицательные примеры и правила допуска в продуктивный контур.
Ранее возможность представления учебных планов с помощью GNN и LSTM также исследовалась авторами в отдельной постановке задачи \cite{Kozhanov2025}; в настоящей работе эти архитектуры повторно проверяются на новом frozen-корпусе ЕПВО и не получают автоматического преимущества перед текстовым baseline.

\subsection{Многокритериальная оптимизация}

Проектирование плана является многокритериальной задачей: нужно максимизировать покрытие LO и междисциплинарность, одновременно минимизируя дублирование, нарушения квот и сложность пререквизитов. В рабочем контуре \systemname{} эту задачу решает детерминированный constrained-планировщик с repair-эвристиками, exact-fit балансированием и независимым верификатором. NSGA-II, строящий Парето-множество посредством недоминируемой сортировки и crowding distance \cite{Deb2002}, реализован отдельно как экспериментальная ветка и не используется для представленных production-планов с явным направлением или группой ЕПВО.

\subsection{Воспроизводимость и эксперт в контуре}

Datasheets for Datasets предлагают документировать происхождение, состав и ограничения набора данных \cite{Gebru2021}, а Model Cards --- назначение модели, процедуру оценки и границы применимости \cite{Mitchell2019}. Интерактивное машинное обучение подчёркивает, что пользователь должен участвовать не только после обучения, но и в формировании обратной связи \cite{Amershi2014}. Поэтому прогноз в \systemname{} маркируется как ``прогноз ИИ'', а эксперт может подтвердить, ослабить, отклонить или исправить связь.

Таким образом, существующие направления решают отдельные части задачи: OBE задаёт проверяемые результаты, recommender/retrieval-подходы находят содержательно близкие элементы, curricular analytics анализирует уже существующий граф, а многокритериальная оптимизация ищет компромиссные решения. Исследовательский разрыв состоит в отсутствии единого контура, который одновременно использует национальную экспертную память, ограничивает кандидатов уровнем и направлением, строит полный кредитно-сбалансированный план, проверяет пререквизиты и сохраняет объяснимое экспертное решение. Именно эту композицию и её воспроизводимые admission-инварианты исследует \systemname{}.

\section{Данные ЕПВО}

\subsection{Источник и содержание карточек}

Исходные данные получены из публичного интерфейса Реестра образовательных программ ЕПВО \cite{EPVORegistry}. Для каждой доступной карточки сохранялась исходная JSON-структура без изменения. В зависимости от заполненности карточка могла содержать:
\begin{itemize}
  \item идентификатор, статус, наименование и цель программы на русском, казахском и английском языках;
  \item уровень образования, область образования, направление подготовки и группу образовательных программ;
  \item общий объём кредитов и продолжительность;
  \item список результатов обучения на трёх языках;
  \item дисциплины, их названия, краткие описания, кредиты, курс и семестр;
  \item заявленные связи дисциплин с результатами обучения;
  \item экспертный уровень связи, интерпретированный как $0$, $0{,}5$ или $1$ для отклонённой, средней и высокой связи;
  \item служебные поля происхождения и идентификаторы исходных объектов.
\end{itemize}

Не каждая карточка содержит все поля, поэтому отсутствие значения не заменялось вымышленной информацией. Raw-слой сохраняет исходный объект и SHA-256 checksum. Это позволяет повторно выполнить нормализацию и проверить, что исходная запись не была незаметно изменена.

\subsection{Слоистая модель хранения}

Данные разделены на четыре уровня:
\begin{enumerate}
  \item \textbf{raw EPVO}: неизменяемые исходные программы, дисциплины, LO и экспертные проверки;
  \item \textbf{normalized EPVO}: очищенные и дедуплицированные сущности, справочники направлений и групп;
  \item \textbf{approved course repository}: дисциплины, разрешённые для рабочего генератора;
  \item \textbf{generation candidates}: ограниченный контекст конкретного проекта с учётом выбранных направлений, групп, языка и квот.
\end{enumerate}

\begin{figure}[H]
\centering
\placeholderfigure{fig02\_epvo\_data\_pipeline.pdf}{Показать четыре слоя данных, provenance, checksum, дедупликацию, трёхъязычные поля, expert labels и переход только approved-записей в генератор.}
\caption{Конвейер подготовки данных ЕПВО с provenance, дедупликацией и контролем качества переводов.}
\label{fig:data-pipeline}
\end{figure}

Нормализация названия выполняется приведением к нижнему регистру, удалением лишних символов и пробелов. Отпечаток дисциплины строится по трём языковым названиям:
\begin{equation}
  f(c)=\operatorname{SHA256}\left(
  n(t_c^{ru})\,\Vert\,n(t_c^{kk})\,\Vert\,n(t_c^{en})
  \right),
  \label{eq:fingerprint}
\end{equation}
где $n(\cdot)$ --- функция нормализации, а $\Vert$ --- конкатенация. При совпадении отпечатка объединяются исходные программы, направления, группы и source keys. Исходные записи при этом остаются в raw-слое.

\subsection{Паспорт набора данных}

\begin{table}[H]
\centering
\caption{Состав набора данных ЕПВО}
\label{tab:dataset}
\begin{tabularx}{\linewidth}{>{\raggedright\arraybackslash}X r}
\toprule
\textbf{Объект} & \textbf{Количество}\\
\midrule
Исходные программы & 12\,215\\
Исходные записи дисциплин & 408\,638\\
Исходные результаты обучения & 124\,521\\
Исходные экспертные проверки / пары & 935\,151\\
Направления подготовки & 125\\
Группы образовательных программ & 483\\
Нормализованные дисциплины & 191\,292\\
Нормализованные экспертные связи & 932\,483\\
Пары с численной экспертной меткой & 800\,793\\
Пары без численной метки & 134\,358\\
\bottomrule
\end{tabularx}
\end{table}

Файл \texttt{course\_lo\_pairs.jsonl} имеет размер 3\,742\,779\,112 байт и SHA-256
\texttt{37a996be8ae9cc1f6c5ce5588612b1e0142e6238a1b856ad0d9afae45845d3a4}.
Файл \texttt{programs.jsonl} имеет размер 23\,876\,141 байт и SHA-256
\texttt{137c5fdcc2ed371789c9e4a4c274e3d5707d63d9d2dd724655d0f6e18b9eb6e4}.

\subsection{Разбиение без утечки программного контекста}

Разбиение выполнялось не по отдельным парам, а по идентификатору программы:
\begin{equation}
  b(p)=
  \operatorname{int}\left(
  \operatorname{SHA256}(42\Vert id_p)_{1:8},16
  \right)\bmod 100.
  \label{eq:split}
\end{equation}
Если $b<70$, программа относится к train; при $70\le b<85$ --- к validation; иначе --- к test. Получены 654\,470, 138\,443 и 142\,238 пар соответственно. Все дисциплины, LO и связи одной программы попадают только в одну часть. Такой подход не устраняет все возможные повторы одинаковых дисциплин между разными университетами, но предотвращает наиболее прямую утечку контекста одной карточки.

\section{Постановка задачи}

Пусть $\Ccal=\{c_1,\ldots,c_n\}$ --- множество дисциплин-кандидатов, $\Ocal=\{o_1,\ldots,o_m\}$ --- результаты обучения, $S=\{1,\ldots,T\}$ --- семестры, а $\Ecal_{pre}\subseteq\Ccal\times\Ccal$ --- отношение пререквизитов. Для дисциплины известны кредиты $cr_i$, рекомендуемый семестр $r_i$, область $d_i$ и текстовое описание. Для пары $(c_i,o_j)$ вычисляется сила связи $w_{ij}\in[0,1]$.

Двоичная переменная $x_i$ показывает включение дисциплины в план, а $z_{is}$ --- назначение дисциплины семестру:
\begin{equation}
  x_i\in\{0,1\},\qquad z_{is}\in\{0,1\},\qquad
  \sum_{s=1}^{T} z_{is}=x_i.
  \label{eq:variables}
\end{equation}

Для программы задаются целевой объём $CR^\ast$, допустимая нагрузка $L_s^{min},L_s^{max}$ и минимальные доли направлений $q_d$. Основные жёсткие ограничения:
\begin{align}
  \sum_i cr_i x_i &= CR^\ast, \label{eq:credits}\\
  L_s^{min}\le \sum_i cr_i z_{is} &\le L_s^{max}, \quad s=1,\ldots,T, \label{eq:load}\\
  z_{is}=1,\ (c_k,c_i)\in\Ecal_{pre}
  &\Rightarrow \sum_{u<s}z_{ku}=1, \label{eq:prerequisite}\\
  \frac{\sum_{i:d_i=d}cr_i x_i}{\sum_i cr_i x_i} &\ge q_d. \label{eq:quota}
\end{align}

Дополнительно запрещаются повторяющиеся нормализованные названия, чужие направления, недопустимые типы компонентов и практика в раннем семестре. Задача состоит в поиске нескольких различных допустимых планов, максимизирующих качество.

\begin{plainbox}
Переменная $x_i$ отвечает на вопрос ``включить ли дисциплину'', а $z_{is}$ ---
``в какой семестр её поставить''. Формула~\eqref{eq:credits} требует точную сумму
кредитов; \eqref{eq:load} не разрешает слишком пустой или перегруженный семестр;
\eqref{eq:prerequisite} ставит основы раньше продолжения; \eqref{eq:quota}
сохраняет заданную долю направления. Поэтому высокий AI-score означает только
смысловую близость и не даёт права нарушить учебные или нормативные правила.
\end{plainbox}

\section{Метод \systemname{}}

\subsection{Мастер параметров программы и scoped repository}

До поиска дисциплин пользователь задаёт:
\begin{itemize}
  \item уровень образования;
  \item классический или междисциплинарный тип программы;
  \item область образования, направление подготовки и группу ОП;
  \item для междисциплинарной программы --- второе отличное направление и группу;
  \item язык обучения, продолжительность, число семестров и общий объём кредитов;
  \item минимальные доли двух направлений;
  \item цель и результаты обучения.
\end{itemize}

Кандидаты извлекаются иерархически:
\begin{equation}
  \text{same group}
  \succ
  \text{same direction}
  \succ
  \text{near group}
  \succ
  \text{evidence-constrained bridge}.
  \label{eq:scope}
\end{equation}
Для классической программы активно одно направление. Для междисциплинарной --- два разных направления и заданная пользователем пропорция. Domain guard запрещает автоматически расширять план случайными дисциплинами из нерелевантных областей.

Репозиторий реализует управляемое продвижение данных между слоями
\begin{equation}
  \mathrm{raw}_{EPVO}
  \rightarrow
  \mathrm{normalized}_{EPVO}
  \rightarrow
  \mathrm{approved}\ repository
  \rightarrow
  \mathrm{generation}\ candidates.
  \label{eq:repository-layers}
\end{equation}
Индексация выполняется для выбранного направления или группы ОП, а связь канонической дисциплины с нормализованной записью хранится явно. Это предотвращает загрузку всех 191\,292 записей непосредственно в активный каталог и сохраняет происхождение данных. Карточка дисциплины содержит отдельные поля \texttt{title\_ru/title\_kk/title\_en} и \texttt{description\_ru/description\_kk/description\_en} со статусом перевода. В контрольной проверке направление магистратуры 7M083 содержало 141 нормализованную запись: 102 новые канонические карточки были созданы, 39 связаны с существующими. Для направления докторантуры 8D051 обработана 221 запись: 193 карточки созданы и 28 связаны с существующими.

Целевой объём программы $C^{target}$ и технический допуск $\delta_C$ задаются явно. Допустимый итог удовлетворяет
\begin{equation}
  C^{target}\le \sum_i cr_i x_i \le C^{target}+\delta_C,
  \qquad \delta_C=3\ \text{по умолчанию}.
  \label{eq:credit-tolerance}
\end{equation}
Планировщик стремится к точному значению $C^{target}$; допуск используется только при дискретном балансировании кредитов и отображается пользователю до генерации.

\subsection{Гибридный поиск}

Текст дисциплины строится из названия, описания, области и доступных метаданных. Результат обучения используется как запрос. SBERT-кодировщик $g_\theta$ формирует нормированные векторы:
\begin{equation}
  \mathbf{u}_i=g_\theta(c_i),\qquad
  \mathbf{v}_j=g_\theta(o_j),\qquad
  s_{ij}^{sem}=
  \frac{\mathbf{u}_i^\top\mathbf{v}_j}
       {\|\mathbf{u}_i\|_2\|\mathbf{v}_j\|_2}.
  \label{eq:cosine}
\end{equation}

Лексический score рассчитывается BM25 \cite{Robertson2009} и нормируется относительно максимума среди кандидатов. Итоговый retrieval score:
\begin{equation}
  r_{ij}=0.75\,s_{ij}^{sem}+0.25\,s_{ij}^{BM25}.
  \label{eq:hybrid}
\end{equation}
Смысловая компонента помогает находить перефразированные соответствия, а BM25 поддерживает точные профессиональные термины, аббревиатуры и названия технологий.

\begin{plainbox}
Если LO содержит ``собирать и сохранять цифровые доказательства'', поиск должен
найти курс ``Цифровая криминалистика'', даже когда формулировки различаются.
SBERT отвечает за близость смысла, а BM25 --- за точные слова вроде Kafka,
MRI или уголовно-процессуального кодекса. Их смесь уменьшает риск пропустить
как перефразирование, так и узкий термин.
\end{plainbox}

\subsection{Оценка связи дисциплина--LO}

В ресурсоэкономном режиме используется гибридная эвристика:
\begin{equation}
  h_{ij}=\clip_{[0,1]}
  \left(
    s_{ij}^{sem}
    +b_{ij}^{kw}
    +b_{ij}^{dom}
    +b_{ij}^{overlap}
  \right),
  \label{eq:heuristic}
\end{equation}
где $b^{kw}\le0.30$ --- бонус ключевых слов, $b^{dom}\le0.15$ --- совпадение области, $b^{overlap}\le0.20$ --- нормированное пересечение слов.

Если локальная SBERT-модель включена, косинусная близость преобразуется в калиброванную уверенность вокруг порога, выбранного на validation:
\begin{equation}
  p_{ij}^{AI}=
  \sigmoid\left(
    \frac{s_{ij}^{sem}-\tau_{AI}}{T_{AI}}
  \right)
  =
  \frac{1}
       {1+\exp\left[-(s_{ij}^{sem}-\tau_{AI})/T_{AI}\right]}.
  \label{eq:calibration}
\end{equation}
Этот score является прогнозом, а не экспертным вердиктом.

Для коротких профессиональных названий SBERT может занижать связь даже при буквальном совпадении термина с LO. Поэтому используется ограниченная лексическая поддержка
\begin{equation}
  b_{ij}^{lex}=\min\!\left(0.20,
  0.70b_{ij}^{kw}+0.50b_{ij}^{overlap}+0.20b_{ij}^{dom}+b_{ij}^{title}\right),
  \qquad
  p_{ij}^{KAG}=\clip_{[0,1]}\left(p_{ij}^{AI}+b_{ij}^{lex}\right),
  \label{eq:ai-lexical-support}
\end{equation}
где $b^{title}\le0.15$ активируется только при совпадении не менее половины содержательных терминов названия дисциплины с LO. Ограничение $0.20$ не позволяет одному лексическому совпадению превратить нерелевантную дисциплину в сильную связь.

\subsection{Перенос экспертного сигнала ЕПВО}

Экспертная оценка в ЕПВО относится к исходному LO конкретной программы. Нельзя механически переносить её на новый LO только потому, что совпала дисциплина. Поэтому используется консервативный текстовый мост. Пусть $K(o)$ --- множество ключевых терминов результата. Сходство нового и исходного LO:
\begin{equation}
  a(o_j,o_q^{src})=
  \frac{|K(o_j)\cap K(o_q^{src})|}
       {\max\left(1,\min(|K(o_j)|,|K(o_q^{src})|)\right)}.
  \label{eq:expert-overlap}
\end{equation}
Экспертный сигнал учитывается только при $a\ge0.12$:
\begin{equation}
  e_{ij}=
  \ell_{iq}^{exp}\cdot
  \min(1,1.5a),
  \label{eq:expert-transfer}
\end{equation}
где $\ell^{exp}\in\{0,0.5,1\}$ --- уровень исходной экспертной оценки. Из нескольких исходных LO выбирается максимальный допустимый сигнал.

Итоговая оценка:
\begin{equation}
  w_{ij}=
  \max\left(
    h_{ij}^{\ast},
    \min(1,0.65h_{ij}^{\ast}+0.35e_{ij})
  \right),
  \label{eq:expert-blend}
\end{equation}
где $h_{ij}^{\ast}=p_{ij}^{KAG}$ при активной модели и $h_{ij}^{\ast}=h_{ij}$ в fallback-режиме. Формула позволяет сильной реальной экспертной оценке поднять уверенность, но не заменяет смысловую проверку.

\begin{plainbox}
В интерфейсе AI-score означает прогноз модели для новой пары
``дисциплина--LO''. EPVO-score означает перенесённое свидетельство из ранее
проверенных программ, скорректированное по сходству формулировок LO. Это не два
процента покрытия программы. Это два разных источника уверенности в одной
связи; окончательное решение остаётся за правилами и экспертом.
\end{plainbox}

\subsection{Экспертная обратная связь}

После генерации эксперт выбирает один из вердиктов: \texttt{confirmed}, \texttt{weak}, \texttt{incorrect}, \texttt{corrected}. Обновлённый score:
\begin{equation}
\tilde w_{ij}=
\begin{cases}
\max(w_{ij},0.75), & \texttt{confirmed},\\
\min(w_{ij},0.45), & \texttt{weak},\\
0.05, & \texttt{incorrect},\\
w_{ij}^{corr}, & \texttt{corrected}.
\end{cases}
\label{eq:feedback}
\end{equation}
Если эксперт указал собственное численное значение, оно имеет приоритет. Отклонённые связи исключаются из блока ``Почему выбрана?'' и из подтверждённых результатов семестра.

\subsection{Вероятностное покрытие результатов обучения}

Максимальный score одной дисциплины не отражает накопление нескольких независимых подтверждений. Поэтому для LO $o_j$ используется вероятностное объединение:
\begin{equation}
  C_j=1-\prod_{i:x_i=1}(1-\tilde w_{ij}).
  \label{eq:prob-coverage}
\end{equation}
Например, две независимые связи силой 0,6 дают
$1-(1-0,6)^2=0,84$, а не 1,2. Однако агрегированное значение не должно скрывать
отсутствие одной сильной реальной дисциплины, поэтому отдельно сохраняются
$R_j=\max_i \tilde w_{ij}$ по реальным курсам, число доказательств и признак
bridge-поддержки.

\begin{plainbox}
Покрытие LO показывает, насколько план в целом подтверждает результат.
Строка ``LO3: 63\%, EPVO 21\%, AI'' в старом интерфейсе могла быть понята
неверно. В новой интерпретации 63\% --- итоговая сила конкретной связи после
объединения сигналов, 21\% --- отдельное экспертное свидетельство ЕПВО, а
``AI'' --- происхождение прогноза. Для оценки всего LO система показывает
список всех дисциплин-источников и отдельно отмечает реальные курсы и bridge.
\end{plainbox}
Если одну компетенцию поддерживают несколько релевантных дисциплин, $C_j$ увеличивается, но остаётся в $[0,1]$. Для отображения подтверждённых связей используется порог 0,40; порог качества полного LO задаётся конфигурацией и в международном чек-листе обычно равен 0,60. Bridge-модуль инициируется, если:
\begin{equation}
  g_j=\max(0,\tau_C-C_j)>0.
  \label{eq:gap}
\end{equation}

\subsection{Граф знаний и пререквизиты}

Граф программы:
\begin{equation}
  \Gcal=(V,E_{pre}\cup E_{LO}\cup E_{sim}),
  \label{eq:graph}
\end{equation}
где $V$ включает дисциплины и LO; $E_{pre}$ --- пререквизиты, $E_{LO}$ --- связи покрытия, $E_{sim}$ --- смысловая близость дисциплин. Для дисциплин ЕПВО консервативные пререквизиты выводятся внутри выбранного направления из типичной последовательности семестров и предметной близости; явно заданные связи имеют приоритет. Фундаментальная дисциплина может иметь экспертную отметку \texttt{prerequisite exempt}; тогда входящие пререквизиты очищаются, а автоматический алгоритм не создаёт для неё искусственную зависимость. Перед оценкой решения выполняется prerequisite closure: включение продвинутой дисциплины добавляет необходимые базовые дисциплины.

\subsection{Многокритериальная оптимизация}

Для выбранного множества $S_x=\{i:x_i=1\}$ оптимизируются следующие показатели.

\paragraph{Ограниченное среднее покрытие}
\begin{equation}
  f_{cov}(x)=\frac{1}{m}
  \sum_{j=1}^{m}
  \min\left(1,\sum_{i\in S_x}\tilde w_{ij}\right).
  \label{eq:objective-coverage}
\end{equation}

\paragraph{Среднее семантическое дублирование}
\begin{equation}
  f_{red}(x)=
  \frac{1}{|S_x|(|S_x|-1)}
  \sum_{\substack{i,k\in S_x\\i\ne k}}
  \cos(\mathbf{u}_i,\mathbf{u}_k).
  \label{eq:redundancy}
\end{equation}

\paragraph{Энтропия направлений}
\begin{equation}
  f_{ent}(x)=
  -\sum_{d\in\Dcal}p_d(x)\ln p_d(x),
  \quad
  p_d(x)=
  \frac{\sum_{i:d_i=d}cr_i x_i}
       {\sum_i cr_i x_i}.
  \label{eq:entropy}
\end{equation}

\paragraph{Штраф квоты}
\begin{equation}
  f_{quota}(x)=
  \sum_{d\in\Dcal}\max(0,q_d-p_d(x)).
  \label{eq:quota-penalty}
\end{equation}

\paragraph{Штраф кредитов и нагрузка пререквизитов}
\begin{align}
  f_{credit}(x)&=
  \begin{cases}
  |CR^\ast-CR(x)|, & CR(x)<CR^\ast,\\
  \max(0,CR(x)-CR^{max}), & CR(x)\ge CR^\ast,
  \end{cases}\\
  f_{pre}(x)&=\sum_{i\in S_x}|\operatorname{Prereq}(c_i)|.
  \label{eq:other-penalties}
\end{align}

Вектор оптимизации:
\begin{equation}
  F(x)=
  \left(
  f_{cov},
  -f_{red},
  -f_{quota},
  f_{ent},
  -f_{credit},
  -f_{pre}
  \right).
  \label{eq:objective-vector}
\end{equation}

Популяция решений проходит prerequisite repair, credit repair, недоминируемую сортировку, crossover, mutation и отбор по crowding distance \cite{Deb2002}. Из различных допустимых Парето-решений выбираются три варианта:
\begin{itemize}
  \item \textbf{A}: приоритет покрытия LO, затем низкого дублирования и квот;
  \item \textbf{B}: приоритет точной кредитной допустимости, затем покрытия и низкого дублирования;
  \item \textbf{C}: приоритет низкого дублирования, квот и покрытия.
\end{itemize}
После выбора одного решения оно удаляется из множества кандидатов, поэтому A/B/C не являются тремя ссылками на один и тот же genome.

\begin{figure}[H]
\centering
\placeholderfigure{fig03\_nsga2\_pareto.pdf}{Показать облако кандидатов и Pareto-front по coverage, redundancy и quota penalty; отметить разные точки A, B, C и последующую жёсткую верификацию.}
\caption{Экспериментальная Парето-ветка и выбор альтернатив A/B/C перед независимой проверкой.}
\label{fig:pareto}
\end{figure}

\subsection{Планировщик семестров и верификатор}

Выбранные дисциплины распределяются по семестрам с учётом рекомендуемого семестра, типа дисциплины, пререквизитов и целевой нагрузки. Практики и итоговая аттестация ограничиваются поздними семестрами. Финальный repair может перемещать дисциплины, заменять или удалять целые элективные единицы, изменять только проектные bridge-модули в явном диапазоне 3--7 кредитов и устранять дубли названий. Кредитность реальной утверждённой дисциплины ЕПВО считается атомарной и автоматически не переписывается.

Поскольку одна каноническая дисциплина может иметь несколько исторических EPVO-карточек,
планировщик перед размещением агрегирует их \texttt{typical\_semester} медианой. Это делает
порядок устойчивым к случайному дублю конкретной организации и применяется только к
реальным EPVO-дисциплинам; нормативные компоненты ГОСО и bridge-модули сохраняют
явно заданные семестры.

\begin{plainbox}
Repair --- это не новое обучение модели, а исправление уже собранного плана.
Например, он может перенести сложный курс из первого семестра в третий или
заменить лишний электив. При этом система не имеет права превратить реальный
пяти-кредитный курс в трёхкредитный только ради удобной арифметики.
\end{plainbox}

Независимый верификатор повторно рассчитывает:
\begin{itemize}
  \item общий объём кредитов;
  \item нагрузку каждого семестра;
  \item отсутствие пропущенного или не более раннего пререквизита;
  \item минимальное и среднее покрытие LO;
  \item число сильных доказательств;
  \item среднее семантическое дублирование.
\end{itemize}
План считается feasible только при нуле жёстких нарушений. Старые планы удаляются или заменяются только после успешного построения всех трёх новых вариантов, что делает генерацию транзакционной с точки зрения пользователя.

\subsection{Bridge-модули и замена реальными дисциплинами}

Bridge-модуль создаётся только при недостаточном покрытии профессионального LO или при необходимости явной связи двух направлений. Сначала система ищет реальную дисциплину:
\begin{equation}
  \text{same EPVO scope}
  \land
  \text{not duplicated}
  \land
  (e_{ij}\ge0.50\ \lor\ p_{ij}^{AI}\ge0.65).
  \label{eq:bridge-replacement}
\end{equation}
Domain/content guard дополнительно проверяет содержание. Например, медицинский bridge не может быть заменён общей дисциплиной по алгоритмам только из-за слабого текстового совпадения. Если доказательной реальной замены нет, bridge сохраняется и явно маркируется как спроектированный модуль, а не как запись ЕПВО.

\subsection{Объяснимость}

Для каждой дисциплины формируется explanation bundle:
\begin{itemize}
  \item наиболее сильные LO и полные тексты результатов;
  \item semantic/heuristic/AI score;
  \item экспертный score ЕПВО и исходная программа;
  \item совпавшие ключевые слова и фрагмент описания;
  \item направление и группа ЕПВО;
  \item пререквизиты и роль в графе;
  \item причина выбора семестра;
  \item вердикт локального эксперта.
\end{itemize}
Таким образом, система отвечает не только на вопрос ``что включить?'', но и на вопрос ``почему это включено?''.

\section{Программная реализация}

\systemname{} реализована как локальное web-приложение с backend API и frontend-интерфейсом. Данные и модели не требуется отправлять во внешний сервис. Основные функциональные модули:
\begin{enumerate}
  \item мастер создания программы;
  \item репозиторий дисциплин с фильтрами направления и группы;
  \item генератор связей course--LO;
  \item детерминированный constrained-планировщик A/B/C; NSGA-II — отдельная экспериментальная ветка;
  \item граф пре- и постреквизитов;
  \item анализ покрытия и достижимости LO;
  \item bridge generator и preview замены;
  \item международный чек-лист OBE/ABET-style/CDIO/Tuning;
  \item генератор тематического плана дисциплины;
  \item экран ``Сравнить с ЕПВО'';
  \item Dataset Passport и Model Benchmark Dashboard;
  \item трёхъязычный интерфейс и отдельные статусы draft/verified для переводов.
\end{enumerate}

Долгая генерация была оптимизирована без увеличения RAM: устранена повторная полная индексация, добавлен LRU-кэш на 8\,192 эмбеддинга, SQLite переведён в WAL, \texttt{busy\_timeout} увеличен до 120 секунд, а параллельное удаление проекта блокируется во время построения. Progress API хранит этап, процент, число обработанных LO, количество найденных связей и время последнего обновления.

LLM-функции не имеют права напрямую менять учебный план. Ответы предложений целей, LO и анализа достижимости проходят единый типизированный контракт Pydantic: проверяются количество элементов, допустимые статусы, диапазоны оценок и упоминание выбранных предметных областей. Если внешний провайдер возвращает нерелевантный или некорректный JSON, система сохраняет детерминированный результат и явно отмечает источник ответа. Этот слой подготовлен для необязательной интеграции PydanticAI без изменения правил ГОСО, EPVO и планировщика.

\section{Экспериментальная методика}

\subsection{Оборудование и воспроизводимость}

Длительные эксперименты выполнялись локально на NVIDIA GeForce RTX 4050 Laptop GPU с CUDA. Зафиксированы:
\begin{itemize}
  \item seed 42;
  \item программно-уровневый split;
  \item локальные пути моделей без повторного скачивания;
  \item JSON-файлы метрик;
  \item команды запуска и журналы;
  \item SHA-256 исходных ML-файлов;
  \item правило допуска новой модели.
\end{itemize}

Новая модель могла заменить рабочую только если на frozen split одновременно улучшались ROC-AUC и PR-AUC, а F1 не снижался. Для ranking-модели требовалось не ухудшать Recall@K, MRR и nDCG.

\subsection{Метрики классификации}

Для бинарной задачи ``существует ли связь'':
\begin{align}
  Precision &= \frac{TP}{TP+FP},\\
  Recall &= \frac{TP}{TP+FN},\\
  F1 &= 2\frac{Precision\cdot Recall}{Precision+Recall}.
  \label{eq:classification-metrics}
\end{align}
ROC-AUC оценивает способность ранжировать положительные примеры выше отрицательных при изменении порога. PR-AUC особенно информативна при дисбалансе классов. Порог классификации выбирался только на validation и затем без изменения применялся к test.

\subsection{Метрики ранжирования}

Для каждого LO формируется список дисциплин-кандидатов. Использовались:
\begin{align}
  Recall@K &=
  \frac{|Rel\cap TopK|}{|Rel|},\\
  MRR &= \frac{1}{Q}\sum_{q=1}^{Q}\frac{1}{rank_q},\\
  DCG@K &= \sum_{r=1}^{K}\frac{2^{rel_r}-1}{\log_2(r+1)},\\
  nDCG@K &= \frac{DCG@K}{IDCG@K}.
  \label{eq:ranking-metrics}
\end{align}
nDCG учитывает позицию и градуированную релевантность \cite{Jarvelin2002}. Ranking benchmark использовал детерминированно выбранные 80 test-программ и 810 LO-запросов.

\subsection{Модели}

\paragraph{Base SBERT.}
\texttt{paraphrase-multilingual-mpnet-base-v2} без предметного дообучения.

\paragraph{SBERT 40k.}
Дообучение на 40\,000 пар ЕПВО, одна эпоха, mixed precision, локальная CUDA. Сохранённый журнал фиксирует итоговый train loss 0,2045 и время 2\,292 секунды. Нижняя часть энкодера была частично заморожена для снижения памяти и риска разрушения многоязычного пространства.

\paragraph{Expert head.}
Небольшая обучаемая голова поверх исходных SBERT-признаков для прогноза численного экспертного уровня.

\paragraph{GNN pilot.}
One-hop mean aggregation по графу программы; целевые рёбра удалялись до агрегации. Использовано до 80 программ на split, 10 эпох.

\paragraph{LSTM pilot.}
Последовательность дисциплин упорядочивалась по семестрам. LSTM получала prefix длиной до 8 дисциплин и вектор LO. Использовано до 20 программ на split в фактическом сохранённом прогоне.

\paragraph{Hard-negative variants.}
Проверены лексические hard negatives с долями 75\% и 25\%, semantic hard-negative mining и ансамбль моделей. Hard negatives и выбор пар являются важной частью обучения dense retrieval и bi-encoder моделей \cite{Xiong2021,Thakur2021}, но их способ формирования должен соответствовать целевой задаче.

\paragraph{Ranking-loss pilot.}
Использована Multiple Negatives Ranking Loss: для каждой положительной пары остальные элементы mini-batch становятся отрицательными кандидатами \cite{Henderson2017,Reimers2019}. Эффективная обучающая выборка составила 4\,000 положительных пар, одна эпоха, batch size 4, learning rate $10^{-6}$; обучение заняло около 407 секунд.

\subsection{Разделение производственного и экспериментального контуров}

\begin{table}[H]
\centering
\caption{Алгоритмический протокол экспериментов с полными планами}
\label{tab:planner-protocol}
\begin{tabular}{p{0.28\linewidth}p{0.27\linewidth}p{0.20\linewidth}p{0.14\linewidth}}
\toprule
\textbf{Эксперимент} & \textbf{Planner} & \textbf{Repair/verifier} & \textbf{Seed, runs}\\
\midrule
AgroTech, Киберследователь, IT--медицина & deterministic constrained A/B/C & level/scope guard, LO repair, exact credit check & 42, 3 варианта\\
Магистратура 7M & deterministic constrained A/B/C & admission guard, exact-fit, transactional verifier & 42, 3 варианта\\
Докторантура D094 & deterministic constrained dry-run A & ГОСО block, exact-fit residual, verifier & 42, 1 вариант\\
NSGA-II pilot & evolutionary experimental branch & общий admission contract и verifier & 42, controlled pilot\\
\bottomrule
\end{tabular}
\end{table}

Основная таблица моделей ниже использует единый frozen test из 6\,000 пар только там, где это явно указано. Expert head, GNN и LSTM рассматриваются как controlled pilots на ограниченных подвыборках и не интерпретируются как прямое ранговое сравнение с основным benchmark. Коэффициенты гибридного поиска $0{,}75/0{,}25$ и пороги admission зафиксированы до финальной регрессии; их полномасштабный sensitivity-анализ не подменяется результатами одной конфигурации.

\section{Результаты}

\subsection{RQ1: предметное дообучение SBERT}

\begin{table}[H]
\centering
\caption{Сравнение моделей на frozen test split}
\label{tab:model-comparison}
\small
\begin{tabular}{lrrrrrr}
\toprule
\textbf{Модель} & \textbf{N} & \textbf{ROC-AUC} & \textbf{PR-AUC} & \textbf{F1} & \textbf{Prec.} & \textbf{Recall}\\
\midrule
Base multilingual SBERT & 6\,000 & 0.7031 & 0.6946 & 0.6886 & 0.5668 & 0.8773\\
SBERT pilot & 6\,000 & 0.7325 & 0.7278 & 0.7039 & --- & ---\\
SBERT 40k (production baseline) & 6\,000 & 0.7666 & 0.7681 & 0.7188 & 0.6255 & 0.8450\\
\textbf{Cached SBERT 12k candidate} & 6\,000 & \textbf{0.7748} & \textbf{0.7746} & \textbf{0.7268} & \textbf{0.6522} & ---\\
Expert head & 1\,674 & 0.6832 & 0.6471 & 0.6581 & 0.6208 & 0.7001\\
GNN one-hop pilot & 2\,082 & 0.5391 & 0.5471 & 0.6684 & --- & ---\\
LSTM sequence pilot & 4\,310 & 0.4911 & 0.5034 & 0.6657 & --- & ---\\
\bottomrule
\end{tabular}
\end{table}

Дообучение SBERT 40k дало абсолютный прирост ROC-AUC 0,0635, PR-AUC 0,0735 и F1 0,0302 относительно base model. Precision вырос на 0,0587, а recall снизился на 0,0323. Cached-модель 12k дополнительно улучшила ROC-AUC на 0,0082, PR-AUC на 0,0065, F1 на 0,0080 и precision на 0,0267 относительно SBERT 40k. Её accuracy составила 0,6915 против 0,6695. Поскольку Recall@5 снизился относительно отдельного ranking-loss pilot, 12k трактуется как лучший кандидат для классификации наличия связи, но не как доказанно лучший ранжировщик и не заменяет production baseline без отдельной валидации.

Результат не равен ``76,66\% правильных ответов''. ROC-AUC измеряет качество порядка пар при всех порогах. Фактическая accuracy при выбранном validation-пороге 0,3449 составила 0,6695. Для практической системы важен не один процент, а сочетание ranking-метрик, экспертных подтверждений и жёсткой проверки плана.

\subsection{RQ2: hard negatives, ансамбль и ranking loss}

\begin{table}[H]
\centering
\caption{Эксперименты по улучшению SBERT}
\label{tab:sbert-ablation}
\small
\begin{tabularx}{\linewidth}{>{\raggedright\arraybackslash}X rrrrrr}
\toprule
\textbf{Вариант} & \textbf{ROC} & \textbf{PR} & \textbf{F1} & \textbf{R@10} & \textbf{MRR} & \textbf{nDCG@10}\\
\midrule
SBERT 40k & 0.7666 & 0.7681 & 0.7188 & 0.7041 & 0.7758 & 0.6882\\
Lexical hard negative 75\% & 0.7552 & 0.7590 & 0.7034 & --- & --- & ---\\
Lexical hard negative 25\% & 0.7583 & 0.7625 & 0.7117 & --- & --- & ---\\
Ensemble 95/5 & \textbf{0.7692} & 0.7676 & \textbf{0.7233} & 0.7033 & 0.7731 & 0.6874\\
Semantic hard negative & 0.7647 & 0.7638 & 0.7177 & 0.7038 & 0.7734 & 0.6864\\
Ranking-loss pilot & 0.7631 & 0.7643 & 0.7133 & \textbf{0.7071} & 0.7815 & \textbf{0.6923}\\
Scoped expert-memory v2 & --- & --- & --- & 0.6855 & \textbf{0.8032} & 0.6808\\
\bottomrule
\end{tabularx}
\end{table}

Лексические hard negatives ухудшили основной baseline. Вероятная причина --- выбор ``трудности'' по пересечению слов не всегда соответствует экспертной смысловой границе. Semantic hard-negative mining восстановил часть качества, но также не превзошёл SBERT 40k.

Ансамбль 95\% SBERT 40k + 5\% balanced hard-negative model немного улучшил ROC-AUC и F1, но ухудшил все измеренные ranking-метрики. Он не был внедрён, поскольку генератору важен порядок дисциплин в top-$K$, а не только бинарное решение.

Ranking-loss pilot дал небольшой, но согласованный прирост относительно SBERT 40k:
\begin{itemize}
  \item Recall@5: $0.5131\rightarrow0.5185$;
  \item Recall@10: $0.7041\rightarrow0.7071$;
  \item MRR: $0.7758\rightarrow0.7815$;
  \item nDCG@10: $0.6882\rightarrow0.6923$.
\end{itemize}
При этом классификационные показатели снизились. Результат поддерживает двухэтапную архитектуру: SBERT 40k определяет вероятность существования связи, а ranking-loss model пересортировывает уже допустимый shortlist.

Дополнительный эксперимент scoped expert-memory v2 использовал только
train-программы: 8\,568 программ, 62\,425 уникальных scoped-свидетельств и
точную группу ЕПВО. На frozen test он улучшил MRR до 0,8032, но снизил
Recall@10 до 0,6855; production-модель поэтому не изменялась. Hit@10 при этом
равнялся 0,9720. Эти числа отвечают на разные вопросы: Hit@10 означает, что
\emph{хотя бы одна} правильная дисциплина встретилась в первой десятке, тогда
как Recall@10 измеряет, какая доля \emph{всех} правильных дисциплин попала в
десятку. Поэтому Hit@10 выше 0,90 не является достижением Recall@10 выше 0,90.

\begin{plainbox}
Если для LO существует десять подходящих дисциплин, а система нашла в top-10
семь, Recall@10 равен 0,7. Но Hit@10 уже равен 1, потому что найден хотя бы один
правильный ответ. Цель Recall@10 $>0,90$ остаётся исследовательской; текущий
подтверждённый максимум 0,7071 нельзя честно называть 90-процентным recall.
\end{plainbox}

\begin{figure}[H]
\centering
\placeholderfigure{fig04\_two\_stage\_model.pdf}{Слева широкий candidate set, затем SBERT 40k link classifier с threshold, справа ranking-loss reranker; поверх показаны EPVO expert score, domain guard и экспертное подтверждение.}
\caption{Двухэтапная схема определения связи и ранжирования кандидатов.}
\label{fig:two-stage}
\end{figure}

\subsection{RQ3: GNN и LSTM}

GNN pilot достиг ROC-AUC 0,5391 и PR-AUC 0,5471. LSTM pilot получил ROC-AUC 0,4911 и PR-AUC 0,5034, то есть не показал полезного разделения классов. Значения F1 около 0,66 не должны интерпретироваться как успех: при выбранном пороге и структуре выборки модель может часто предсказывать доминирующий положительный класс. ROC-AUC и PR-AUC показывают, что ранжирующая способность остаётся около случайной.

Отрицательный результат научно важен по трём причинам:
\begin{enumerate}
  \item текст дисциплины и LO уже несёт сильный сигнал, который малый графовый pilot не превзошёл;
  \item граф ЕПВО содержит неполные или выведенные пререквизиты, поэтому one-hop aggregation получает шумную структуру;
  \item LSTM требует более точной последовательной постановки, чем короткий prefix и слабая связь ``предыдущие дисциплины--текущий LO''.
\end{enumerate}
Поэтому GNN/LSTM сохранены как воспроизводимые controlled experiments, но не внедрены в генератор. Полный эксперимент имеет смысл после расширения графа, построения hard negatives на уровне программ и оценки на одинаковом внешнем holdout.

\subsection{RQ4: генерация полного плана}

Для прикладной проверки использовался не демонстрационный пример, а набор контрольных пользовательских программ с сохранёнными версиями и независимым повторным аудитом после записи планов в БД. Проверка выполнялась по инвариантам кредита, нагрузки семестра, уровня образования, доменной квоты, порядка пререквизитов, наличия реального подтверждения каждого LO и количества bridge-модулей. Время включает локальный поиск по ЕПВО, SBERT-scoring, построение A/B/C, repair и транзакционное сохранение.

\begin{table}[H]
\centering
\caption{Последняя контрольная регрессия генератора (19 июля 2026 г.)}
\label{tab:plan-case}
\begin{tabular}{p{0.25\linewidth}ccccc}
\toprule
\textbf{Контрольная программа} & \textbf{A/B/C} & \textbf{Кредиты} & \textbf{Hard} & \textbf{Real LO} & \textbf{Bridge}\\
\midrule
ИКТ--агрономия (AgroTech) & 3/3 & 240 & 0 & 9/9 & 3\\
Право--ИКТ (Киберследователь) & 3/3 & 240 & 0 & 9/9 & 1\\
ИКТ--медицина & 3/3 & 240 & 0 & 12/12 & 0\\
Магистратура, 7M, контроль & 3/3 & 120 & 0 & 9/9 & 2\\
Докторантура D094, dry-run A & 1/1 & 180 & 0 & 10/10 & 0\\
\bottomrule
\end{tabular}
\end{table}

В AgroTech все варианты имеют 240 кредитов, ноль hard violations и реальное покрытие 9/9 LO; минимальная эффективная связь равна соответственно 1,0000, 1,0000 и 0,9827. Семестровые нагрузки A равны $27,30,32,32,33,27,31,28$, B --- $28,29,28,33,33,30,31,28$, C --- $27,29,33,31,33,27,30,30$ кредитов. Полный прогон занял 216,08~с.

В программе ``Киберследователь'' исходный аудит обнаружил для A/B/C соответственно $4/9$, $0/9$, $0/9$ LO с реальным подтверждением, $4/9/9$ bridge-модулей и $5/6/5$ дисциплин неверного уровня. После level guard, агрегации доменных свидетельств и финального LO-repair все три варианта получили 240 кредитов, 9/9 real LO, один bridge и ноль дисциплин неверного уровня. Эффективные объёмы направлений составили A: 133,5/106,5, B: 132,5/107,5 и C: 137,5/102,5 кредита; квоты 50\% и 40\% выполнены с допуском 3 кредита.

В программе ИКТ--медицина промежуточная версия имела real LO $11/12$, $11/12$, $12/12$ и по три bridge. После запрета принудительного bridge при полном реальном покрытии и повторного финального repair получены три варианта по 240 кредитов, real LO $12/12$, ноль bridge и ноль hard violations. Полный прогон занял 178,82~с.

После исправления единого admission-контроля constrained-планировщика контрольная магистратура
прошла A/B/C: каждый вариант содержит 120 кредитов, 9/9 реальных LO, ноль
hard violations и 100\% по внутреннему международному checklist. Два bridge
сохранились как явные проектные элементы и не подменяют реальное LO-покрытие.

Для докторантуры создан независимый временный проект направления 8D061,
группы D094. После scoring 800 пар новый точный подбор профильного остатка
сформировал 180 кредитов: нормативный блок ГОСО 155 кредитов и пять реальных
профильных дисциплин по 5 кредитов. Получены real LO 10/10, bridge 0,
hard violations 0, международный checklist 100\% и нагрузки
$30/30/30/30/27/33$. Технический проект после проверки удалён, а JSON-отчёт
сохранён. Это regression-проверка алгоритмических инвариантов, а не решение
аккредитационного органа и не оценка будущей успеваемости студентов.

Высокое вероятностное покрытие не означает, что каждый bridge должен автоматически считаться окончательной дисциплиной. Большое число bridge является диагностикой недостатка реальных дисциплин в выбранном scope или слишком широких LO. Поэтому интерфейс показывает предупреждение, preview реальных замен и источники покрытия ``real course vs bridge''.

Baseline-аудит существующей базы выявил 12 проектов и 36 планов. Число неразличимых A/B/C составило 0. Последняя регрессия дополнительно проверяет сохранённые планы повторным чтением из БД: результат генерации считается успешным только тогда, когда все три варианта одновременно проходят одинаковый независимый верификатор.

\subsection{RQ5: внешнее структурное сравнение с программами ЕПВО}

Чтобы не ограничиваться самопроверкой собственными правилами, семь сгенерированных программ были сопоставлены с полными учебными планами ЕПВО той же группы или направления. Для каждого плана вычислялись: доля реальных дисциплин с происхождением из ЕПВО, согласованность семестра с типичным семестром реестра в пределах $\pm1$, максимальный Jaccard с отдельной референтной программой и доля сгенерированного плана, присутствующая в ближайшей программе. Сравнение выполнялось по каноническим идентификаторам дисциплин, seed~42; для средних рассчитан непараметрический bootstrap по программам, 10,000 повторов.

\begin{table}[H]
\centering
\caption{Внешнее структурное сравнение с полными программами ЕПВО}
\label{tab:external-epvo-validation}
\begin{tabular}{p{0.31\linewidth}cccc}
\toprule
\textbf{Программа} & \textbf{EPVO provenance} & \textbf{Semester $\pm1$} & \textbf{Best Jaccard} & \textbf{Containment}\\
\midrule
AgroTech & 1,000 & 0,809 & 0,054 & 0,064\\
Киберследователь & 0,958 & 0,739 & 0,046 & 0,065\\
IT--медицина & 1,000 & 0,643 & 0,069 & 0,143\\
Интеллектуальное госуправление & 0,400 & 0,833 & 0,087 & 0,333\\
Цифровой агроном & 0,480 & 0,083 & 0,143 & 0,167\\
Цифровые двойники и роботы & 0,729 & 0,429 & 0,041 & 0,086\\
Киберзащита инфраструктуры & 0,740 & 0,270 & 0,100 & 0,135\\
\midrule
Среднее [bootstrap 95\% CI] & 0,758 [0,581; 0,917] & 0,510 [0,329; 0,681] & 0,077 [0,054; 0,103] & 0,142 [0,088; 0,214]\\
Среднее шести quality-eligible планов & 0,805 & 0,581 & 0,066 & 0,138\\
\bottomrule
\end{tabular}
\end{table}

Результат показывает, что 75,8\% реальных дисциплин всех контрольных планов имеют прямое происхождение из реестра ЕПВО, а 51,0\% дисциплин с известным типовым семестром размещены в пределах одного семестра от медианного семестра подходящих EPVO-источников. Использование медианы по канонической дисциплине предотвращает зависимость результата от случайного дубля одной программы. Защищённые регуляторные элементы ГОСО (практики, исследовательская работа и аттестация) не имеют обязательного EPVO-course-id и поэтому считаются отдельным источником, а не ошибкой происхождения: после их исключения EPVO provenance quality-eligible планов составляет 99,3\%. Один намеренно созданный отрицательный stress-test с невыполнимой нагрузкой не включается в production-среднее; для шести quality-eligible планов общая семестровая согласованность составляет 58,1\%. Низкий Jaccard с одной ближайшей программой ожидаем для новых междисциплинарных профессий: генератор не копирует один существующий план, а собирает подтверждённые элементы нескольких программ выбранных направлений. Поэтому overlap интерпретируется как внешний структурный ориентир и проверка provenance, а не как доказательство педагогической эффективности. Слепая содержательная экспертиза независимыми специалистами остаётся отдельным уровнем валидации.

Дополнительный транзакционный прогон после внедрения scoped-медианы пересобрал семь
A-вариантов без изменения сохранённых пользовательских планов. В свежем прогоне
EPVO-provenance без нормативных компонентов составил 0,9964, а согласование семестра
— 0,553. Историческое значение 0,581 относится к планам, созданным до изменения
планировщика; оба значения сохранены, чтобы не смешивать регрессионную совместимость
и эффект новой версии алгоритма.

  \subsection{Сопоставление с готовыми рекомендованными планами ЕПВО}

Для внешней структурной проверки использовались готовые полные планы ЕПВО той же
группы или направления, а не только внутренние правила системы. Сравнение выполнялось
по трём осям: происхождение дисциплины из подходящего EPVO-скоупа, согласование типового
семестра с медианой аналогичных карточек и структурное пересечение с ближайшей полной
программой. Для шести выполнимых контрольных программ среднее EPVO-provenance составило
0,8046, а после исключения защищённых нормативных компонентов ГОСО — 0,993; доля
дисциплин в пределах одного семестра от медианного типового семестра составила 0,581.
Один намеренно невыполнимый stress-test отделён от quality-eligible выборки и не входит
в производственное среднее. Такой тест проверяет происхождение и структурную правдоподобность
плана, но не заменяет слепую содержательную оценку независимыми экспертами.

\subsection{RQ6: объяснимость и участие эксперта}

В интерфейсе эксперт видит полный LO при наведении, численные AI/EPVO scores, источник дисциплины, причину семестра и пререквизиты. Нелогичная связь, например клинической дисциплины с общекомандным LO только по слабому совпадению, может быть помечена как \texttt{incorrect}; после этого она получает score 0,05 и исключается из подтверждённого покрытия. Такая обратная связь формирует новый локальный обучающий слой.

\begin{figure}[H]
\centering
\placeholderfigure{fig05\_human\_feedback\_loop.pdf}{Цикл: prediction -> explanation -> expert confirmed/weak/incorrect/corrected -> feedback store -> recalculation -> future retraining. Показать разделение «прогноз ИИ» и «экспертное решение».}
\caption{Контур экспертной обратной связи и накопления проверенных решений.}
\label{fig:feedback-loop}
\end{figure}

\section{Обсуждение}

\subsection{Что было извлечено из экспертной базы}

Работа с ЕПВО дала три типа знания:
\begin{enumerate}
  \item \textbf{каталожное знание}: какие дисциплины реально используются в конкретных направлениях и группах;
  \item \textbf{структурное знание}: типичные кредиты, семестры, принадлежность направлениям и повторяемость дисциплин между программами;
  \item \textbf{экспертное знание}: какие дисциплины связывались с какими LO и с какой силой.
\end{enumerate}

Это отличает систему от генератора, использующего только общую языковую модель. Название дисциплины не считается достаточным доказательством. Решение опирается на текст, реестр, экспертную связь, направление, граф и ограничения плана.

\subsection{Научная новизна}

Новизна работы заключается не в изобретении SBERT, BM25, NSGA-II или графов по отдельности, а в их новой воспроизводимой композиции для национального контекста проектирования образовательных программ:
\begin{enumerate}
  \item впервые в рамках проекта рассеянные экспертные оценки ЕПВО преобразованы в контролируемый ML-корпус с программно-уровневым split и checksum;
  \item предложена двойная доказательная оценка course--LO, объединяющая предметно дообученную семантическую модель и переносимый экспертный сигнал;
  \item предложено вероятностное накопление покрытия LO с явным разделением реальных дисциплин и bridge-модулей;
  \item генерация сформулирована как constrained multi-objective optimization, а не как свободное текстовое продолжение;
  \item обоснована двухэтапная схема ``link classifier + ranking reranker'' на основании разных направлений улучшения метрик;
  \item создан human-in-the-loop цикл, в котором ошибка интерфейса становится структурированной меткой для последующего обучения;
  \item обеспечена трёхъязычная provenance-aware модель дисциплин со статусами draft/verified;
  \item предложена агрегация доменных свидетельств по всем исходным строкам канонической дисциплины, устраняющая зависимость результата от порядка записей;
  \item сформулирован единый admission-контракт для эвристического и NSGA-II путей: каждая реальная дисциплина должна сохранять аудируемые LO-, level- и scope-свидетельства до финальной записи плана;
  \item предложен exact-fit алгоритм профильного остатка после ГОСО, совместно оптимизирующий целые кредиты реальных дисциплин и покрытие профессиональных LO;
  \item введено двухпороговое разделение реального покрытия и bridge-гипотезы, а также лексикографический repair, минимизирующий bridge только после выполнения hard constraints;
  \item реализован транзакционный протокол A/B/C с повторной верификацией сохранённого состояния БД и сохранением старого активного плана при любой ошибке.
\end{enumerate}

Нельзя без систематического обзора заявлять, что система абсолютно единственная в мире. Корректная формулировка состоит в том, что \systemname{} представляет оригинальную интеграцию экспертных данных национального реестра, KAG-поиска, формального планировщика и explainable feedback в одном работающем прототипе.

\subsection{Практическая ценность}

Система полезна:
\begin{itemize}
  \item разработчикам новых и междисциплинарных образовательных программ;
  \item академическим комитетам при сравнении вариантов учебного плана;
  \item экспертам реестра при поиске слабых или неподтверждённых связей;
  \item службам качества при проверке кредитов, LO coverage и пререквизитов;
  \item исследователям образовательных данных как воспроизводимый benchmark.
\end{itemize}

Система не принимает решение об утверждении программы. Она сокращает пространство вариантов, выявляет нарушения и предоставляет доказательства, после чего эксперт подтверждает или корректирует результат.

\section{Дополнительная верификация и угрозы валидности}

\begin{enumerate}
  \item \textbf{Неполные отрицательные метки.} Отсутствие связи в карточке не всегда означает доказанную нерелевантность.
  \item \textbf{Исторический bias.} Экспертные решения отражают практику существующих программ и могут закреплять устаревшие структуры.
  \item \textbf{Повторы между учреждениями.} Program-level split предотвращает прямую утечку карточки, но одинаковая дисциплина может встречаться в разных программах.
  \item \textbf{Ограниченный ranking benchmark.} Использованы 80 программ и 810 запросов; внешнее структурное сравнение с ЕПВО и bootstrap по семи контрольным программам добавлены, однако необходим отдельный institutional/temporal holdout для ML-метрик.
  \item \textbf{GNN/LSTM являются pilot, а не финальным сравнением.} Их размер и постановка пока недостаточны для вывода о классе архитектур в целом.
  \item \textbf{Двухэтапная схема ещё требует production validation.} Ranking-loss улучшил порядок, но не должен заменять классификатор.
  \item \textbf{Пререквизиты частично выведены.} Не все карточки ЕПВО содержат явные связи, поэтому часть графа строится консервативно по семестрам и смысловой близости.
  \item \textbf{Переводы имеют статус качества.} Автоматически созданное описание не должно считаться verified до проверки.
  \item \textbf{Высокое вероятностное покрытие может насыщаться.} Несколько средних связей дают высокий $C_j$, поэтому дополнительно учитываются число сильных доказательств и максимальная одиночная связь.
  \item \textbf{Нет проспективного исследования результатов студентов.} Оценено качество структуры и экспертных связей, но не причинное влияние плана на успеваемость выпускников.
  \item \textbf{Докторский контроль пока ограничен одним dry-run вариантом A.} Он подтверждает исправление exact-fit и инвариантов D094, но не заменяет многогрупповую проверку A/B/C и внешнюю экспертную оценку содержания дисциплин.
\end{enumerate}

\subsection{Проверенные инварианты constrained curriculum synthesis}

Последняя версия метода разделяет \emph{предсказание смысловой связи} и
\emph{допустимость включения дисциплины}. Высокий score не отменяет жёсткие
ограничения уровня, направления, кредита и пререквизитов. Это ключевое отличие
от генерации плана одним запросом к языковой модели.

\subsection{Контроль образовательного уровня и аудируемая включаемость дисциплины}

В рабочей системе добавлено жёсткое ограничение уровня образования по кодам
ЕПВО: $6B/B$ для бакалавриата, $7M/M$ для магистратуры и $8D/D$ для
докторантуры. Для дисциплины $c$ и проекта $p$ индикатор допустимости уровня
определяется как
\begin{equation}
I_{level}(c,p)=\mathbb{I}\left[\exists d\in D_c:d\sim prefix(level_p)\;\lor\;
\exists g\in G_c:g\sim groupPrefix(level_p)\right].
\end{equation}
Итоговое правило допуска профессиональной дисциплины имеет вид
\begin{equation}
I_{admit}(c,p)=I_{scope}(c,p)I_{level}(c,p)
\mathbb{I}\left[\max(s_{AI},s_{EPVO})\geq\tau\right],\qquad \tau=0.4.
\label{eq:admission}
\end{equation}
Таким образом, текстовое совпадение предметной области не может перенести курс
магистратуры в бакалавриат. Дисциплина без прямой связи с хотя бы одним
результатом обучения считается жёстким нарушением. Для обязательных компонентов
ГОСО РК автоматически создаются отдельные системные результаты обучения уровня
программы; это делает нормативные дисциплины объяснимыми в OBE-матрице, а не
формальными исключениями. Эксперт может пакетно подтвердить несколько замен
bridge-модулей, причём решение сохраняется и воспроизводится при следующем
построении A/B/C. Аналитическая панель сравнивает варианты, дедуплицирует пары
``дисциплина--РО'' и показывает измеренное усиление каждого результата.

\subsection{Агрегация доменных свидетельств после дедупликации}

Одна каноническая дисциплина $c$ может соответствовать множеству исходных
записей $\mathcal{R}(c)$ из разных программ ЕПВО. Ранее последовательное чтение
этих строк могло ошибочно оставить домен последней записи. Предложена
детерминированная агрегация максимального свидетельства:
\begin{equation}
E_{c,d}=\max_{r\in\mathcal{R}(c)} scope(r,d),\qquad
scope(r,d)=
\begin{cases}
3,& group(r)=group(d),\\
2,& direction(r)=direction(d),\\
0,& \text{иначе},
\end{cases}
\label{eq:domain-evidence}
\end{equation}
\begin{equation}
d^{\ast}(c)=\arg\max_{d\in\mathcal{D}_p}E_{c,d},
\label{eq:domain-assignment}
\end{equation}
где при равенстве используется заранее заданный приоритет основного
направления. Тем самым принадлежность определяется всем происхождением
канонической дисциплины, а не порядком строк в БД.

Эффективный объём направления вычисляется как
\begin{equation}
C_d(P)=\sum_{c\in P}cr_c\,\mathbb{I}[d^{\ast}(c)=d]+B_d(P),
\label{eq:effective-domain-credits}
\end{equation}
где $B_d(P)$ --- только явно определённая доля предметных и интеграционных
bridge-модулей. Общий LO-gap bridge не считается доказательством домена.
Ограничение квоты имеет вид
\begin{equation}
C_d(P)+\delta_d\geq \rho_d CR^{\ast},
\label{eq:domain-quota-tolerance}
\end{equation}
где $\rho_d$ --- заданная доля, $\delta_d$ --- кредитный допуск.

\subsection{Точный подбор профильного остатка после ГОСО РК}

В программах, создаваемых в режиме ГОСО РК, нормативный блок практик,
исследовательской работы и итоговой аттестации добавляется отдельно. Если его
объём равен $C^{GOSO}$, для реальных профильных дисциплин остаётся
\begin{equation}
  C^{prof}=C^{target}-C^{GOSO}.
  \label{eq:professional-residual}
\end{equation}
Обычное жадное удаление курсов может оставить неудобную сумму и затем заполнить
разницу bridge-модулями. Поэтому введено ограниченное динамическое
программирование. Для кандидата $c_i$ определяется битовая маска $m_i$
профессиональных LO, подтверждённых с силой не менее 0,50, и полезность
$u_i=\sum_{\ell\in m_i}s_{i\ell}^{eff}$. Состояние
\begin{equation}
  DP[q,m]=\max \sum_i u_i x_i
  \quad\text{при}\quad
  \sum_i cr_i x_i=q,
  \quad \bigvee_{i:x_i=1}m_i=m
  \label{eq:goso-knapsack}
\end{equation}
строится только для $0\le q\le C^{prof}$. Среди состояний с
$q=C^{prof}$ лексикографически выбирается максимальное число покрытых LO, затем
максимальная доказательность и, при равенстве, меньшее число дисциплин.

\begin{plainbox}
Для докторантуры из 180 кредитов нормативные элементы заняли 155, поэтому
осталось ровно 25 профильных кредитов. Новый алгоритм не берёт первые попавшиеся
курсы и не меняет их кредитность: он ищет комбинацию реальных дисциплин ровно на
25 кредитов, которая закрывает как можно больше профессиональных результатов.
В контрольном D094 это дало пять реальных курсов по 5 кредитов и ноль bridge.
\end{plainbox}

\subsection{Двухпороговое реальное покрытие и минимизация bridge}

Для пары дисциплина--LO используется эффективное свидетельство
\begin{equation}
s_{c\ell}^{eff}=\max\{s_{c\ell}^{SBERT},s_{c\ell}^{EPVO}\}.
\label{eq:effective-lo-link}
\end{equation}
Реальное покрытие не смешивается с проектным bridge:
\begin{equation}
R_{\ell}(P)=\max_{c\in P:\,real(c)=1}s_{c\ell}^{eff}.
\label{eq:real-lo-coverage}
\end{equation}
Порог $\tau_r=0.50$ используется при repair для поиска реальной дисциплины,
а $\tau_q=0.60$ --- для строгого отчёта качества. Значение 0,75 у bridge
является проектной гипотезой и показывается отдельно; оно не может превратить
bridge в экспертно подтверждённую дисциплину.

Число bridge минимизируется лексикографически после выполнения hard constraints:
\begin{equation}
\min_P\left(
N_{hard}(P),\;N_{miss}^{real}(P),\;N_{bridge}(P),\;
\Pi_{quota}(P),\;\Pi_{load}(P),\;-Q(P)
\right),
\label{eq:lexicographic-repair}
\end{equation}
где $N_{miss}^{real}=\sum_{\ell}\mathbb{I}[R_{\ell}<\tau_r]$.
Если $N_{miss}^{real}=0$ и bridge отсутствуют, алгоритм запрещает искусственно
добавлять ``улучшающий'' bridge. После каждого квотного, диверсификационного или
экспертного repair все LO проверяются повторно, что устраняет потерю покрытия
на поздних стадиях.

Кандидат $c$ для замены bridge $b$ ранжируется функцией
\begin{equation}
J(c\mid b)=w_1|M_b\cap L_c|+w_2\mathbb{I}[s^{EPVO}>0]
+w_3\max_{\ell\in M_b}s_{c\ell}^{eff}
+w_4 rank_{scope}(c)-w_5|cr_c-cr_b|,
\label{eq:bridge-candidate-ranking}
\end{equation}
где $M_b$ --- LO bridge. Подтверждённая экспертом замена становится
неизменяемым ограничением следующей генерации.

\subsection{Семестровая последовательность и транзакционная проверка A/B/C}

Нагрузка семестра
\begin{equation}
L_s(P)=\sum_c cr_c z_{cs},\qquad
L_s^{target}-\Delta_s\le L_s(P)\le L_s^{target}+\Delta_s,
\label{eq:semester-load-new}
\end{equation}
проверяется совместно с $z_{pre(c),s'}<z_{c,s}$ для $s'<s$. Семантические
окна запрещают преждевременное размещение продвинутых потоковых, forensic,
клинических и исследовательских дисциплин, если необходимые основы ещё не
пройдены. Repair может перераспределять 3--7-кредитные проектные модули, но не
нарушать итоговый допуск и сохранённые экспертные решения.

Различие вариантов измеряется симметрической разностью
\begin{equation}
D(P_i,P_j)=|\mathcal{C}(P_i)\triangle\mathcal{C}(P_j)|,
\qquad i\ne j.
\label{eq:variant-diversity}
\end{equation}
Новый набор публикуется атомарно только при
\begin{equation}
Commit(A,B,C)=\mathbb{I}\left[
\bigwedge_{v\in\{A,B,C\}} Feasible(v)\land
\bigwedge_{i<j}D(P_i,P_j)>0\right].
\label{eq:transactional-commit}
\end{equation}
При $Commit=0$ старый активный план сохраняется. После commit планы повторно
читаются из БД и проходят тот же независимый verifier; поэтому проверяется не
только объект в памяти, но и фактически сохранённый результат.

\subsection{Что именно проверено}

Метод подтверждён regression-тестами на трёх различных междисциплинарных
бакалаврских программах, контрольной магистратуре и независимом dry-run
докторантуры. Во всех девяти сохранённых бакалаврских вариантах выполнены
целевые кредиты, нулевое число hard violations и реальное покрытие всех LO;
магистерские A/B/C повторили эти инварианты на 120 кредитах; докторский вариант
прошёл их на 180 кредитах без bridge. Измеренные изменения образуют малую
абляцию: level guard снизил число курсов неверного уровня в
``Киберследователе'' с $5/6/5$ до нуля; финальный LO-repair поднял real LO до
$9/9$; правило no-forced-bridge снизило bridge в ИКТ--медицине с трёх до нуля;
exact-fit после ГОСО снизил bridge в D094 с пяти до нуля и поднял
профессиональное real LO с $0/5$ до $5/5$. Следовательно, проверены программные
инварианты и воспроизводимость генерации, но не причинное влияние плана на
результаты студентов и не решение аккредитационного органа.

\section{Направления дальнейшего исследования}

\begin{enumerate}
  \item Провести bootstrap-анализ ROC-AUC, PR-AUC, Recall@K, MRR и nDCG с доверительными интервалами.
  \item Разделить внешний holdout по университетам, направлениям и времени регистрации программы.
  \item Ввести calibrated probability и reliability diagrams отдельно по уровням 6B, 7M и 8D.
  \item Расширить пройденную регрессию магистратуры и докторантуры на несколько независимых групп 7M/8D и все варианты A/B/C докторантуры.
  \item Обучать expert head на накопленной локальной обратной связи confirmed/weak/incorrect/corrected.
  \item Повторить GNN и LSTM только после очистки полного графа и построения программно-зависимых hard negatives.
  \item Провести blinded evaluation с несколькими независимыми экспертами, измерить Cohen/Fleiss $\kappa$ и сравнить время ручного проектирования с \systemname{}.
\end{enumerate}

\section{Критерии готовности и следующие десять шагов}

Готовность оценивалась раздельно, чтобы не смешивать программную корректность,
качество модели и педагогическую валидность. На текущем этапе функциональная стабильность
и регрессионные проверки составляют 90\%, целостность EPVO-репозитория и локализация
RU/KK/EN — 99\%, качество классификации связей — около 80\% от практического
production-порога, структурное соответствие готовым EPVO-планам — 80\% по provenance
и 58\% по согласованию семестра $\pm1$, объяснимость и аудит — 85\%. Внешняя
педагогическая валидность пока оценивается в 50\%, поскольку blinded-оценка
независимыми экспертами ещё не проведена. Это инженерный статус, а не заявление
об аккредитации или доказательство улучшения результатов студентов.

Наиболее полезные следующие шаги: (1) blinded-оценка нескольких программ;
(2) temporal/institutional holdout; (3) bootstrap-доверительные интервалы и
калибровка вероятностей; (4) sensitivity-анализ весов и порогов; (5) инкрементальное
обновление EPVO без полной переиндексации; (6) автоматическое восстановление
пререквизитов с подтверждением спорных рёбер; (7) генерация силлабусов и тематических
планов; (8) расширение магистерских и докторских контрольных программ; (9) единый
контейнер воспроизводимости; (10) сравнение с ручным проектированием по времени,
числу исправлений и экспертному качеству. Эти шаги переводят систему от внутренне
проверяемого прототипа к внешне подтверждённому инструменту поддержки решений.

\subsection{Почему PydanticAI является опциональным слоем}

PydanticAI полезен как адаптер внешней LLM: он задаёт типизированную схему ответа для
предложений целей, результатов обучения и анализа достижимости, проверяет обязательные
поля и позволяет безопасно повторить запрос или перейти к детерминированному fallback.
Это уменьшает риск ошибки ``\texttt{[object Object]}'', неполного JSON и ответа не на
выбранном языке. В Curriculum-KAG такой слой не выбирает дисциплины и не меняет план:
кредиты, уровень, область, ГОСО, пререквизиты и финальная проверка остаются за
детерминированным планировщиком и независимым верификатором.

Поэтому PydanticAI не является новой моделью качества и не заменяет SBERT или EPVO.
Его роль — контрактный шлюз между LLM и системой ограничений. В production он должен
включаться только после отдельного теста версии провайдера, таймаутов, стоимости,
логирования и fallback; по умолчанию адаптер отключён, что сохраняет воспроизводимость
и не делает работу системы зависимой от внешнего API.

\section{Заключение}

Представлена \systemname{} --- работающая система поддержки проектирования образовательных программ, объединяющая экспертные данные ЕПВО, гибридный поиск, граф знаний, семантическую модель, многокритериальную оптимизацию и экспертную обратную связь. В отличие от свободной генерации, система сначала ограничивает область реальными направлениями и группами, затем рассчитывает доказательные связи дисциплина--LO, строит три альтернативных плана и независимо проверяет кредиты, семестры и пререквизиты.

Предметное дообучение multilingual SBERT на данных ЕПВО существенно улучшило ROC-AUC, PR-AUC и F1. Ranking-loss дал дополнительное улучшение поискового порядка, что обосновывает двухэтапную архитектуру. GNN и LSTM в контролируемых pilot-экспериментах не превзошли текстовый baseline и потому не были внедрены. Такой отрицательный результат является частью научной воспроизводимости: новая архитектура принимается не по сложности или популярности, а только по подтверждённым метрикам.

Последняя проверка показала, что хорошей модели недостаточно, если последующая
логика теряет её кандидатов. Исправление admission-контракта NSGA-II, разделение
профессиональных и нормативных LO, запрет изменения кредитов реальных курсов и
точный knapsack-подбор остатка после ГОСО преобразовали контроль D094 из пяти
bridge и нулевого реального покрытия профессиональных LO в пять реальных
дисциплин, $5/5$ профессиональных LO и ноль bridge при тех же 180 кредитах.
Следовательно, научный вклад включает не только обучение эмбеддингов, но и
формальный способ безопасно превратить вероятностный прогноз в допустимый план.

Главный результат работы --- переход от ``ИИ предлагает список дисциплин'' к проверяемому инженерному процессу:
\[
\text{данные}
\rightarrow
\text{доказательные связи}
\rightarrow
\text{граф}
\rightarrow
\text{оптимизация}
\rightarrow
\text{верификация}
\rightarrow
\text{экспертное решение}.
\]
Это делает систему пригодной для дальнейшей научной оценки, институционального пилотирования и развития в сторону полноценной платформы evidence-based curriculum engineering.

\begin{plainbox}
Главная идея работы проста: ИИ предлагает содержательно подходящие варианты,
математический планировщик собирает их в целое, правила отбрасывают невозможное,
а человек видит доказательства и принимает окончательное решение. Поэтому
система не обещает ``автоматически создать идеальную программу'', но заметно
сокращает пространство ручного поиска и делает каждое решение проверяемым.
\end{plainbox}

\section*{Доступность данных}
Исходные карточки были получены из общедоступного интерфейса ЕПВО и использованы только в объёме, необходимом для исследования образовательных программ; персональные данные не включались в ML-корпус. Для воспроизводимости публикуются схема нормализации, Dataset Passport, контрольные суммы, обезличенный пример пар course--LO и скрипты построения split. Полный raw-слой не перераспространяется автоматически; доступ к нему должен соответствовать условиям источника и применимым правилам управления данными.

\section*{Доступность кода}
Версионированный исходный код, конфигурации экспериментов, проверки качества и инструкция развёртывания размещены в репозитории \url{https://github.com/kmypat-commits/curriculum-kag}. Крупные веса моделей и дамп PostgreSQL публикуются как отдельные артефакты с SHA-256, чтобы репозиторий оставался пригодным для клонирования.

\section*{Этика и управление данными}
Исследование не использовало данные студентов, оценки успеваемости, медицинские сведения или иные персональные данные и не вмешивалось в учебный процесс. ЕПВО применялось как источник описаний программ, дисциплин, LO и обезличенных экспертных связей. Результаты системы являются поддержкой решения: окончательное утверждение программы остаётся за уполномоченными экспертами и организацией образования.

\section*{Финансирование}
Исследование выполнено без отдельного внешнего финансирования, если иное не будет указано в метаданных итоговой публикации.

\section*{Конфликт интересов}
Автор заявляет об отсутствии конфликта интересов.

\section*{Вклад авторов}
Концептуализация, постановка задачи, разработка программной системы, подготовка данных, проведение экспериментов, анализ результатов и подготовка рукописи выполнены автором; конкретное распределение ролей CRediT уточняется при добавлении соавторов.

\section*{Раскрытие использования генеративного ИИ}
Генеративный ИИ использовался как инструмент программной разработки, языкового редактирования и подготовки вариантов формулировок. Все алгоритмы, численные результаты, ссылки и выводы проверялись автором; ИИ не выступает автором и не принимает экспертных решений об утверждении образовательных программ.

\iffalse
\appendix

\section{Какие скриншоты вставить в статью}
\label{app:screenshots}

Ниже перечислены интерфейсные иллюстрации. Перед публикацией необходимо скрыть персональные данные и идентификаторы, не относящиеся к исследованию. Рекомендуемая ширина исходного изображения --- не менее 1600 пикселей.

\begin{longtable}{p{0.06\linewidth}p{0.27\linewidth}p{0.58\linewidth}}
\toprule
\textbf{№} & \textbf{Экран} & \textbf{Что должно быть видно}\\
\midrule
\endfirsthead
\toprule
\textbf{№} & \textbf{Экран} & \textbf{Что должно быть видно}\\
\midrule
\endhead
S1 & Мастер создания ОП, шаг 1 &
Уровень образования, тип ``Междисциплинарная'', первое и второе направление ЕПВО, группы ОП и квоты. Подпись: ``Ограничение пространства генерации направлениями ЕПВО''.\\

S2 & Цель и результаты обучения &
Три предложенные цели или LO, возможность выбрать и отредактировать. При наведении на LO должен отображаться полный текст.\\

S3 & Plan Builder A/B/C &
Три действительно разных варианта, активная кнопка выбора любого варианта, 240 кредитов и progress/status.\\

S4 & Карточка ``Почему выбрана?'' &
Корректная профессиональная дисциплина, LO, AI score, EPVO score, источник, причина семестра и кнопки экспертной обратной связи. Не использовать заведомо ошибочную связь как положительный пример.\\

S5 & Граф пререквизитов &
Слева дисциплины семестра, в центре граф связей, справа ``что получает студент после семестра''. Показать пре- и постреквизиты.\\

S6 & LO Coverage &
Покрытие каждого LO с разделением ``реальные дисциплины'', ``bridge'' и ``weak'', а также полное описание LO.\\

S7 & Bridge replacement &
Предупреждение о большом числе bridge, реальные кандидаты ЕПВО, AI/EPVO scores и кнопку ``Добавить в проект''.\\

S8 & Международный чек-лист &
OBE/ABET-style/CDIO/Tuning, общий score, выявленные нарушения и кнопку автоматического исправления последовательности/кредитов.\\

S9 & Сравнить с ЕПВО &
Похожие программы, типовые дисциплины, отсутствующие элементы, уникальность и рекомендации.\\

S10 & Dataset Passport &
Количество программ, дисциплин, LO, expert links, split, seed и checksum.\\

S11 & Model Benchmark &
Таблица SBERT base/40k/expert head/GNN/LSTM и ranking-метрики.\\

S12 & Тематический план дисциплины &
Название, кредиты, 15-недельный план, методы оценивания, связанные LO и экспорт.\\

S13 & Контроль ГОСО: магистратура и докторантура &
Итоговые кредиты, реальные профильные дисциплины, нормативные практики/НИР/аттестация, нагрузка семестров, 100\% LO и отсутствие bridge в контрольном D094.\\
\bottomrule
\end{longtable}

\section{Промпты для генерации научных иллюстраций}
\label{app:prompts}

\noindent\colorbox{yellow!25}{\parbox{0.96\linewidth}{
\textbf{ПРОМПТ ДЛЯ ГЕНЕРАЦИИ ИЛЛЮСТРАЦИИ 1 --- архитектура системы.}\\
Create a clean publication-grade vector diagram for a scientific paper titled
``Curriculum-KAG''. White background, restrained blue, cyan and violet palette,
no 3D, no decorative AI brains, no photorealism. Show a left-to-right pipeline:
Programme Wizard (education level, EPVO direction 1/2, quotas, goals, learning
outcomes) $\rightarrow$ Scoped EPVO Repository $\rightarrow$ Hybrid Retrieval
(0.75 SBERT cosine + 0.25 BM25) $\rightarrow$ Course--LO Scoring (AI prediction,
EPVO expert score, local expert feedback) $\rightarrow$ Curriculum Knowledge
Graph (courses, prerequisites, LO edges) $\rightarrow$ NSGA-II Multi-objective
Planner $\rightarrow$ A/B/C Plans $\rightarrow$ Hard Verifier. Add a feedback
arrow from Expert Confirmation back to the scoring dataset. Use concise Russian
labels, consistent typography, SVG/PDF-ready, 16:9.
}}

\medskip
\noindent\colorbox{yellow!25}{\parbox{0.96\linewidth}{
\textbf{ПРОМПТ ДЛЯ ГЕНЕРАЦИИ ИЛЛЮСТРАЦИИ 2 --- данные ЕПВО.}\\
Create a scientific data-engineering pipeline diagram on a white background.
Four clearly separated layers: ``Raw EPVO'' with immutable JSON cards and
checksums; ``Normalized EPVO'' with deduplication, canonical titles, RU/KK/EN
fields, directions, groups and expert course--LO links; ``Approved Course
Repository'' with verified status; ``Generation Candidates'' filtered by the
selected programme scope. Show provenance arrows and a side branch to
``Dataset Passport: counts, split, seed, SHA-256''. Use flat vector design,
professional journal style, no futuristic glow, Russian labels, 16:9.
}}

\medskip
\noindent\colorbox{yellow!25}{\parbox{0.96\linewidth}{
\textbf{ПРОМПТ ДЛЯ ГЕНЕРАЦИИ ИЛЛЮСТРАЦИИ 3 --- Парето-оптимизация.}\\
Draw a publication-quality conceptual Pareto-front chart for curriculum
optimization. Axes: ``LO coverage (maximize)'', ``Semantic redundancy
(minimize)'', with point colour representing ``Domain quota penalty''. Show a
cloud of candidate plans, a highlighted Pareto front, and three distinct points
A, B, C with small callouts: A coverage-first, B credit-feasibility-first,
C diversity-first. To the right show a hard verifier gate checking 240 credits,
semester load and prerequisite order. Minimal vector style, white background,
Russian labels, suitable for an academic article.
}}

\medskip
\noindent\colorbox{yellow!25}{\parbox{0.96\linewidth}{
\textbf{ПРОМПТ ДЛЯ ГЕНЕРАЦИИ ИЛЛЮСТРАЦИИ 4 --- двухэтапная модель.}\\
Create a clean scientific diagram of a two-stage course-to-learning-outcome
retrieval model. Input: one programme learning outcome and a scoped EPVO course
catalogue. Stage 1: ``SBERT 40k link classifier'' filters candidates using a
validation-selected threshold and EPVO expert evidence. Stage 2:
``Multiple-Negatives ranking reranker'' sorts the shortlist. Final output:
top courses with AI score, EPVO score, evidence and ``requires expert
confirmation''. Include a small metric panel: classifier ROC-AUC 0.7666,
PR-AUC 0.7681; reranker Recall@10 0.7071, MRR 0.7815, nDCG@10 0.6923.
White background, blue/green academic palette, vector, Russian labels.
}}

\medskip
\noindent\colorbox{yellow!25}{\parbox{0.96\linewidth}{
\textbf{ПРОМПТ ДЛЯ ГЕНЕРАЦИИ ИЛЛЮСТРАЦИИ 5 --- human-in-the-loop.}\\
Create an academic circular workflow diagram: ``AI prediction'' $\rightarrow$
``Explanation bundle'' $\rightarrow$ ``Expert decision'' with four branches
Confirmed, Weak, Incorrect, Corrected $\rightarrow$ ``Score recalculation'' and
``Feedback dataset'' $\rightarrow$ ``Future retraining''. Visually distinguish
machine suggestion from human authority. Add a note: ``AI does not approve the
programme''. Flat vector style, white background, accessible colours, Russian
labels, no humanoid robot icons.
}}

\medskip
\noindent\colorbox{yellow!25}{\parbox{0.96\linewidth}{
\textbf{ПРОМПТ ДЛЯ ГЕНЕРАЦИИ ИЛЛЮСТРАЦИИ 6 --- проверенная эволюция качества планов.}\\
Create a publication-grade BEFORE/AFTER scientific infographic, landscape 16:9,
white background, flat vector style, Russian labels, typography similar to
Inter/Helvetica, colour-blind-safe palette. Use three horizontal lanes. Lane 1:
``Право--ИКТ / Киберследователь'', BEFORE cards ``real LO A/B/C = 4/9,
0/9, 0/9; bridge = 4/9/9; wrong education level = 5/6/5'', transformation
blocks ``level guard'', ``aggregated domain evidence'', ``final LO repair'',
AFTER card ``all A/B/C: 240 credits, real LO 9/9, bridge 1, hard violations 0,
wrong level 0''. Lane 2: ``ИКТ--медицина'', BEFORE ``real LO 11/12, 11/12,
12/12; bridge 3/3/3'', transformation ``no forced bridge + repeated repair'',
AFTER ``real LO 12/12, bridge 0, hard 0''. Lane 3: ``ИКТ--агрономия'', AFTER
``A/B/C 240 credits, real LO 9/9, bridge 3, hard 0''. Add a small legend:
green = verified invariant, amber = diagnostic bridge, red = violation removed.
Do not use robots, brains, glowing cyberpunk effects, fake seals or accreditation
logos. Do not invent percentages. Output as editable SVG and 300-dpi PNG.
}}

\medskip
\noindent\colorbox{yellow!25}{\parbox{0.96\linewidth}{
\textbf{ПРОМПТ ДЛЯ ГЕНЕРАЦИИ ИЛЛЮСТРАЦИИ 7 --- агрегация доменных свидетельств.}\\
Draw a precise mathematical data-lineage diagram for an academic journal. On
the left show four immutable EPVO source records that have different programme
groups but normalize to one canonical course. In the centre show a deduplication
node ``canonical course c'' and an evidence matrix with rows source records and
columns Domain 1 and Domain 2. Cell values must use exactly: 3 = exact programme
group, 2 = exact direction, 0 = no evidence. Show the operation
$E_{c,d}=\max_{r\in\mathcal{R}(c)}scope(r,d)$, then deterministic assignment
$d^*(c)=\arg\max_d E_{c,d}$. On the right show domain-credit accounting and the
quota inequality $C_d(P)+\delta_d\ge\rho_d CR^*$. Add a crossed-out mini-panel
``last database row wins'' to explain the eliminated defect. White background,
navy/teal/orange palette, thin lines, no gradients, Russian labels, vector SVG,
A4-width readable at 8-point journal font. Negative prompt: no decorative AI
icons, no fictional data, no 3D database cylinders, no tiny unreadable text.
}}

\medskip
\noindent\colorbox{yellow!25}{\parbox{0.96\linewidth}{
\textbf{ПРОМПТ ДЛЯ ГЕНЕРАЦИИ ИЛЛЮСТРАЦИИ 8 --- транзакционный A/B/C и семестровый граф.}\\
Create a rigorous sequence-and-graph hybrid figure for a scientific paper,
landscape 16:9. Top half: sequence diagram ``User -- Generator -- Scoped EPVO --
SBERT scorer -- Planner -- Independent verifier -- Database''. Show numbered
steps: freeze constraints; retrieve level-safe courses; score course--LO links;
build A, B, C in memory; repair credits/LO/prerequisites; verify every variant;
atomic commit only if all pass; reload from database; verify again. Add an error
branch where any failure preserves the old active plan. Bottom half: an
eight-semester prerequisite DAG, foundation nodes on the left, advanced nodes
later, arrows always forward, semester load bars around 30 credits with a
$\pm3$ tolerance band. Mark real LO evidence with solid green edges and bridge
hypotheses with dashed amber edges. Include formulas
$D(P_i,P_j)=|P_i\triangle P_j|$ and $Commit(A,B,C)=1$ only when all variants are
feasible and distinct. Russian labels, white background, restrained academic
blue/green/amber palette, editable SVG plus 300-dpi PNG. Negative prompt: no
neon glow, no robot, no blockchain symbols, no fake metrics, no backward
prerequisite arrows, no illegible labels.
}}

\medskip
\noindent\colorbox{yellow!25}{\parbox{0.96\linewidth}{
\textbf{ПРОМПТ ДЛЯ ГЕНЕРАЦИИ ИЛЛЮСТРАЦИИ 9 --- точный профильный блок после ГОСО.}\\
Create a publication-grade explanatory vector figure, landscape 16:9, for a
Russian scientific paper. Title: ``Точный подбор профильного блока после ГОСО
РК''. Left panel: a 180-credit doctoral programme represented as one horizontal
stacked bar: 155 credits ``Нормативный блок ГОСО'' and 25 credits
``Профильные дисциплины D094''. Centre panel: show a bounded dynamic-programming
grid with credit states q = 0...25 and five professional LO bits; candidate
course cards must have whole immutable credits, for example five cards of 5
credits. Show the objective: exact q=25, maximize covered professional LO,
then maximize SBERT/EPVO evidence. Right panel: BEFORE card ``5 synthetic
bridge, professional real LO 0/5, domain quota failed''; arrow labelled
``unified NSGA-II admission + exact-fit knapsack + atomic course credits'';
AFTER card ``5 real EPVO courses, bridge 0, professional real LO 5/5,
180/180 credits, semester loads 30/30/30/30/27/33, hard violations 0''.
Distinguish normative LO from professional LO. White background, flat vector,
navy/teal/green palette with amber only for BEFORE bridge, readable Russian
labels, editable SVG and 300-dpi PNG. Negative prompt: no robot, no AI brain,
no neon, no fake university seals, no invented metrics, no 3D, no illegible
equations.
}}

\section{Файлы воспроизводимости проекта}
\label{app:reproducibility}

\begin{itemize}
  \item Dataset Passport: \texttt{backend/experiment-results/dataset-passport.json};
  \item baseline report: \texttt{backend/experiment-results/reproducible-baseline/baseline-report.json};
  \item SBERT 40k benchmark: \texttt{backend/experiment-results/epvo-sbert-finetuned-40k-benchmark/metrics.json};
  \item ranking benchmark: \texttt{backend/experiment-results/epvo-ranking-benchmark/metrics.json};
  \item ranking-loss benchmark: \texttt{backend/experiment-results/epvo-ranking-ranking-loss-pilot/metrics.json};
  \item scoped expert-memory v2: \texttt{backend/experiment-results/epvo-ranking-scoped-memory-v2/metrics.json};
  \item GNN metrics: \texttt{backend/experiment-results/lstm-gnn-controlled-run/gnn-smoke/metrics.json};
  \item LSTM metrics: \texttt{backend/experiment-results/lstm-gnn-controlled-run/lstm-smoke/metrics.json};
  \item article experiment report: \texttt{backend/experiment-results/article-experiment-report/article-experiment-report.json};
  \item контроль AgroTech: \texttt{backend/experiment-results/control-program-audit/report-13-final.json};
  \item контроль ``Киберследователь'': \texttt{backend/experiment-results/control-program-audit/report-15-final-pass.json};
  \item контроль ИКТ--медицина: \texttt{backend/experiment-results/control-program-audit/report-16-complete.json};
  \item контроль докторантуры D094: \texttt{backend/experiment-results/cross-level/doctorate-180-goso-passed-audit.json};
  \item контроль магистратуры после exact-fit: \texttt{backend/experiment-results/cross-level/master-control-after-doctorate-fixes.json};
  \item русское резюме regression-аудита: \texttt{docs/CONTROL\_REGRESSION\_RESULTS\_RU.md}.
\end{itemize}

\fi

\begin{thebibliography}{99}

\bibitem{Biggs1996}
Biggs J.
\emph{Enhancing teaching through constructive alignment}.
Higher Education, 32 (1996), 347--364.
\href{https://doi.org/10.1007/BF00138871}{doi:10.1007/BF00138871}.

\bibitem{ABET2025}
ABET.
\emph{Criteria for Accrediting Computing Programs, 2025--2026}.
\href{https://www.abet.org/accreditation/accreditation-criteria/criteria-for-accrediting-computing-programs-2025-2026/}{Official criteria}.

\bibitem{CDIO2022}
Worldwide CDIO Initiative.
\emph{CDIO Standards 3.0}.
\href{https://www.cdio.org/content/cdio-standards-30}{Official standard}.

\bibitem{Tuning2003}
Gonz\'alez J., Wagenaar R. (eds.).
\emph{Tuning Educational Structures in Europe: Final Report, Phase One}.
University of Deusto and University of Groningen (2003).
\href{https://tuningacademy.org/wp-content/uploads/2014/02/TuningEUI_Final-Report_EN.pdf}{Official report}.

\bibitem{EPVORegistry}
Национальный центр развития высшего образования.
\emph{Реестр образовательных программ высшего и послевузовского образования}.
\href{https://www.enic-kazakhstan.edu.kz/ru/reestr-op/reestr-op-1}{Официальная страница реестра}.

\bibitem{AdiletRegistry}
Министерство науки и высшего образования Республики Казахстан.
\emph{Правила ведения реестра образовательных программ}.
\href{https://adilet.zan.kz/rus/docs/V2200030139}{ИПС ``Әділет'', приказ №106 от 12.10.2022}.

\bibitem{Devlin2019}
Devlin J., Chang M.-W., Lee K., Toutanova K.
\emph{BERT: Pre-training of Deep Bidirectional Transformers for Language Understanding}.
NAACL-HLT (2019), 4171--4186.
\href{https://doi.org/10.18653/v1/N19-1423}{doi:10.18653/v1/N19-1423}.

\bibitem{Reimers2019}
Reimers N., Gurevych I.
\emph{Sentence-BERT: Sentence Embeddings using Siamese BERT-Networks}.
EMNLP-IJCNLP (2019), 3982--3992.
\href{https://doi.org/10.18653/v1/D19-1410}{doi:10.18653/v1/D19-1410}.

\bibitem{Reimers2020}
Reimers N., Gurevych I.
\emph{Making Monolingual Sentence Embeddings Multilingual using Knowledge Distillation}.
EMNLP (2020), 4512--4525.
\href{https://doi.org/10.18653/v1/2020.emnlp-main.365}{doi:10.18653/v1/2020.emnlp-main.365}.

\bibitem{Song2020}
Song K., Tan X., Qin T., Lu J., Liu T.-Y.
\emph{MPNet: Masked and Permuted Pre-training for Language Understanding}.
Advances in Neural Information Processing Systems 33 (2020).
\href{https://proceedings.neurips.cc/paper/2020/hash/c3a690be93aa602ee2dc0ccab5b7b67e-Abstract.html}{NeurIPS paper}.

\bibitem{Thakur2021}
Thakur N., Reimers N., Daxenberger J., Gurevych I.
\emph{Augmented SBERT: Data Augmentation Method for Improving Bi-Encoders for Pairwise Sentence Scoring Tasks}.
NAACL-HLT (2021), 296--310.
\href{https://doi.org/10.18653/v1/2021.naacl-main.28}{doi:10.18653/v1/2021.naacl-main.28}.

\bibitem{Karpukhin2020}
Karpukhin V., Oguz B., Min S., Lewis P., Wu L., Edunov S., Chen D., Yih W.-t.
\emph{Dense Passage Retrieval for Open-Domain Question Answering}.
EMNLP (2020), 6769--6781.
\href{https://doi.org/10.18653/v1/2020.emnlp-main.550}{doi:10.18653/v1/2020.emnlp-main.550}.

\bibitem{Lewis2020}
Lewis P., Perez E., Piktus A., et al.
\emph{Retrieval-Augmented Generation for Knowledge-Intensive NLP Tasks}.
Advances in Neural Information Processing Systems 33 (2020).
\href{https://proceedings.neurips.cc/paper/2020/hash/6b493230-Abstract.html}{NeurIPS paper}.

\bibitem{Liang2024}
Liang L., Sun M., Gui Z., et al.
\emph{KAG: Boosting LLMs in Professional Domains via Knowledge Augmented Generation}.
arXiv:2409.13731 (2024).
\href{https://doi.org/10.48550/arXiv.2409.13731}{doi:10.48550/arXiv.2409.13731}.

\bibitem{Hogan2021}
Hogan A., Blomqvist E., Cochez M., et al.
\emph{Knowledge Graphs}.
ACM Computing Surveys, 54(4), Article 71 (2021).
\href{https://doi.org/10.1145/3447772}{doi:10.1145/3447772}.

\bibitem{Robertson2009}
Robertson S., Zaragoza H.
\emph{The Probabilistic Relevance Framework: BM25 and Beyond}.
Foundations and Trends in Information Retrieval, 3(4) (2009), 333--389.
\href{https://doi.org/10.1561/1500000019}{doi:10.1561/1500000019}.

\bibitem{Deb2002}
Deb K., Pratap A., Agarwal S., Meyarivan T.
\emph{A Fast and Elitist Multiobjective Genetic Algorithm: NSGA-II}.
IEEE Transactions on Evolutionary Computation, 6(2) (2002), 182--197.
\href{https://doi.org/10.1109/4235.996017}{doi:10.1109/4235.996017}.

\bibitem{Kipf2017}
Kipf T.N., Welling M.
\emph{Semi-Supervised Classification with Graph Convolutional Networks}.
ICLR (2017).
\href{https://openreview.net/forum?id=SJU4ayYgl}{OpenReview}.

\bibitem{Hamilton2017}
Hamilton W., Ying Z., Leskovec J.
\emph{Inductive Representation Learning on Large Graphs}.
Advances in Neural Information Processing Systems 30 (2017).
\href{https://proceedings.neurips.cc/paper/2017/hash/5dd9db5e033da9c6fb5ba83c7a7ebea9-Abstract.html}{NeurIPS paper}.

\bibitem{Hochreiter1997}
Hochreiter S., Schmidhuber J.
\emph{Long Short-Term Memory}.
Neural Computation, 9(8) (1997), 1735--1780.
\href{https://doi.org/10.1162/neco.1997.9.8.1735}{doi:10.1162/neco.1997.9.8.1735}.

\bibitem{Xiong2021}
Xiong L., Xiong C., Li Y., et al.
\emph{Approximate Nearest Neighbor Negative Contrastive Learning for Dense Text Retrieval}.
ICLR (2021).
\href{https://openreview.net/forum?id=zeFrfgyZln}{OpenReview}.

\bibitem{Henderson2017}
Henderson M., Al-Rfou R., Strope B., et al.
\emph{Efficient Natural Language Response Suggestion for Smart Reply}.
arXiv:1705.00652 (2017).
\href{https://arxiv.org/abs/1705.00652}{arXiv}.

\bibitem{Jarvelin2002}
J\"arvelin K., Kek\"al\"ainen J.
\emph{Cumulated Gain-based Evaluation of IR Techniques}.
ACM Transactions on Information Systems, 20(4) (2002), 422--446.
\href{https://doi.org/10.1145/582415.582418}{doi:10.1145/582415.582418}.

\bibitem{Heileman2018}
Heileman G.L., Hickman M., Slim A., Abdallah C.T.
\emph{Characterizing the Complexity of Curricula with Directed Graphs}.
arXiv:1811.09676 (2018).
\href{https://arxiv.org/abs/1811.09676}{arXiv}.

\bibitem{Gebru2021}
Gebru T., Morgenstern J., Vecchione B., Vaughan J.W., Wallach H., Daum\'e III H., Crawford K.
\emph{Datasheets for Datasets}.
Communications of the ACM, 64(12) (2021), 86--92.
\href{https://doi.org/10.1145/3458723}{doi:10.1145/3458723}.

\bibitem{Mitchell2019}
Mitchell M., Wu S., Zaldivar A., et al.
\emph{Model Cards for Model Reporting}.
FAT* (2019), 220--229.
\href{https://doi.org/10.1145/3287560.3287596}{doi:10.1145/3287560.3287596}.

\bibitem{Amershi2014}
Amershi S., Cakmak M., Knox W.B., Kulesza T.
\emph{Power to the People: The Role of Humans in Interactive Machine Learning}.
AI Magazine, 35(4) (2014), 105--120.
\href{https://doi.org/10.1609/aimag.v35i4.2513}{doi:10.1609/aimag.v35i4.2513}.

\bibitem{Kozhanov2025}
Kozhanov M., V\'arkonyi-K\'oczy A.R.
\emph{Modelling Curricula with Graph Neural Networks and Long Short-Term Memory Networks}.
Eurasian Journal of Mathematical and Computer Applications, 13(4) (2025), 159--167.

\end{thebibliography}

\end{document}
